% Options for packages loaded elsewhere
\PassOptionsToPackage{unicode}{hyperref}
\PassOptionsToPackage{hyphens}{url}
\documentclass[
]{abntex2}
\usepackage{xcolor}
\usepackage{amsmath,amssymb}
\setcounter{secnumdepth}{-\maxdimen} % remove section numbering
\usepackage{iftex}
\ifPDFTeX
  \usepackage[T1]{fontenc}
  \usepackage[utf8]{inputenc}
  \usepackage{textcomp} % provide euro and other symbols
\else % if luatex or xetex
  \usepackage{unicode-math} % this also loads fontspec
  \defaultfontfeatures{Scale=MatchLowercase}
  \defaultfontfeatures[\rmfamily]{Ligatures=TeX,Scale=1}
\fi
\usepackage{lmodern}
\ifPDFTeX\else
  % xetex/luatex font selection
\fi
% Use upquote if available, for straight quotes in verbatim environments
\IfFileExists{upquote.sty}{\usepackage{upquote}}{}
\IfFileExists{microtype.sty}{% use microtype if available
  \usepackage[]{microtype}
  \UseMicrotypeSet[protrusion]{basicmath} % disable protrusion for tt fonts
}{}
\makeatletter
\@ifundefined{KOMAClassName}{% if non-KOMA class
  \IfFileExists{parskip.sty}{%
    \usepackage{parskip}
  }{% else
    \setlength{\parindent}{0pt}
    \setlength{\parskip}{6pt plus 2pt minus 1pt}}
}{% if KOMA class
  \KOMAoptions{parskip=half}}
\makeatother
\usepackage{color}
\usepackage{fancyvrb}
\newcommand{\VerbBar}{|}
\newcommand{\VERB}{\Verb[commandchars=\\\{\}]}
\DefineVerbatimEnvironment{Highlighting}{Verbatim}{commandchars=\\\{\}}
% Add ',fontsize=\small' for more characters per line
\newenvironment{Shaded}{}{}
\newcommand{\AlertTok}[1]{\textcolor[rgb]{1.00,0.00,0.00}{\textbf{#1}}}
\newcommand{\AnnotationTok}[1]{\textcolor[rgb]{0.38,0.63,0.69}{\textbf{\textit{#1}}}}
\newcommand{\AttributeTok}[1]{\textcolor[rgb]{0.49,0.56,0.16}{#1}}
\newcommand{\BaseNTok}[1]{\textcolor[rgb]{0.25,0.63,0.44}{#1}}
\newcommand{\BuiltInTok}[1]{\textcolor[rgb]{0.00,0.50,0.00}{#1}}
\newcommand{\CharTok}[1]{\textcolor[rgb]{0.25,0.44,0.63}{#1}}
\newcommand{\CommentTok}[1]{\textcolor[rgb]{0.38,0.63,0.69}{\textit{#1}}}
\newcommand{\CommentVarTok}[1]{\textcolor[rgb]{0.38,0.63,0.69}{\textbf{\textit{#1}}}}
\newcommand{\ConstantTok}[1]{\textcolor[rgb]{0.53,0.00,0.00}{#1}}
\newcommand{\ControlFlowTok}[1]{\textcolor[rgb]{0.00,0.44,0.13}{\textbf{#1}}}
\newcommand{\DataTypeTok}[1]{\textcolor[rgb]{0.56,0.13,0.00}{#1}}
\newcommand{\DecValTok}[1]{\textcolor[rgb]{0.25,0.63,0.44}{#1}}
\newcommand{\DocumentationTok}[1]{\textcolor[rgb]{0.73,0.13,0.13}{\textit{#1}}}
\newcommand{\ErrorTok}[1]{\textcolor[rgb]{1.00,0.00,0.00}{\textbf{#1}}}
\newcommand{\ExtensionTok}[1]{#1}
\newcommand{\FloatTok}[1]{\textcolor[rgb]{0.25,0.63,0.44}{#1}}
\newcommand{\FunctionTok}[1]{\textcolor[rgb]{0.02,0.16,0.49}{#1}}
\newcommand{\ImportTok}[1]{\textcolor[rgb]{0.00,0.50,0.00}{\textbf{#1}}}
\newcommand{\InformationTok}[1]{\textcolor[rgb]{0.38,0.63,0.69}{\textbf{\textit{#1}}}}
\newcommand{\KeywordTok}[1]{\textcolor[rgb]{0.00,0.44,0.13}{\textbf{#1}}}
\newcommand{\NormalTok}[1]{#1}
\newcommand{\OperatorTok}[1]{\textcolor[rgb]{0.40,0.40,0.40}{#1}}
\newcommand{\OtherTok}[1]{\textcolor[rgb]{0.00,0.44,0.13}{#1}}
\newcommand{\PreprocessorTok}[1]{\textcolor[rgb]{0.74,0.48,0.00}{#1}}
\newcommand{\RegionMarkerTok}[1]{#1}
\newcommand{\SpecialCharTok}[1]{\textcolor[rgb]{0.25,0.44,0.63}{#1}}
\newcommand{\SpecialStringTok}[1]{\textcolor[rgb]{0.73,0.40,0.53}{#1}}
\newcommand{\StringTok}[1]{\textcolor[rgb]{0.25,0.44,0.63}{#1}}
\newcommand{\VariableTok}[1]{\textcolor[rgb]{0.10,0.09,0.49}{#1}}
\newcommand{\VerbatimStringTok}[1]{\textcolor[rgb]{0.25,0.44,0.63}{#1}}
\newcommand{\WarningTok}[1]{\textcolor[rgb]{0.38,0.63,0.69}{\textbf{\textit{#1}}}}
\usepackage{longtable,booktabs,array}
\usepackage{calc} % for calculating minipage widths
% Correct order of tables after \paragraph or \subparagraph
\usepackage{etoolbox}
\makeatletter
\patchcmd\longtable{\par}{\if@noskipsec\mbox{}\fi\par}{}{}
\makeatother
% Allow footnotes in longtable head/foot
\IfFileExists{footnotehyper.sty}{\usepackage{footnotehyper}}{\usepackage{footnote}}
\makesavenoteenv{longtable}
\usepackage{graphicx}
\makeatletter
\newsavebox\pandoc@box
\newcommand*\pandocbounded[1]{% scales image to fit in text height/width
  \sbox\pandoc@box{#1}%
  \Gscale@div\@tempa{\textheight}{\dimexpr\ht\pandoc@box+\dp\pandoc@box\relax}%
  \Gscale@div\@tempb{\linewidth}{\wd\pandoc@box}%
  \ifdim\@tempb\p@<\@tempa\p@\let\@tempa\@tempb\fi% select the smaller of both
  \ifdim\@tempa\p@<\p@\scalebox{\@tempa}{\usebox\pandoc@box}%
  \else\usebox{\pandoc@box}%
  \fi%
}
% Set default figure placement to htbp
\def\fps@figure{htbp}
\makeatother
\ifLuaTeX
\usepackage[bidi=basic]{babel}
\else
\usepackage[bidi=default]{babel}
\fi
\babelprovide[main,import]{brazilian}
% get rid of language-specific shorthands (see #6817):
\let\LanguageShortHands\languageshorthands
\def\languageshorthands#1{}
\setlength{\emergencystretch}{3em} % prevent overfull lines
\providecommand{\tightlist}{%
  \setlength{\itemsep}{0pt}\setlength{\parskip}{0pt}}
\usepackage[]{biblatex}
\addbibresource{documents\_tcc/refs.bib}
\newcommand{\standalonefiglabel}[1]{\refstepcounter{figure}\label{#1}}
\setcounter{tocdepth}{4}
\usepackage{bookmark}
\IfFileExists{xurl.sty}{\usepackage{xurl}}{} % add URL line breaks if available
\urlstyle{same}
\hypersetup{
  pdftitle={GPU-Accelerated Isoalgorítmico GA+2-opt para TSP},
  pdfauthor={Lucas Galdino},
  pdflang={pt-BR},
  hidelinks,
  pdfcreator={LaTeX via pandoc}}

\title{GPU-Accelerated Isoalgorítmico GA+\(2\)-opt para TSP}
\author{Lucas Galdino}
\date{2025-11-22}

\begin{document}
\maketitle

\clearpage

\section*{Resumo}\label{resumo}
\addcontentsline{toc}{section}{Resumo}

Este trabalho investiga o impacto de diferentes estratégias de
paralelização em Unidades de Processamento Gráfico (GPUs) no desempenho
de um algoritmo memético (Algoritmo Genético híbrido com busca local
2-opt) para o Problema do Caixeiro Viajante (TSP). Foram implementadas e
comparadas quatro variantes isoalgorítmicas: uma versão sequencial em
CPU, uma versão híbrida ingênua (transferência indivíduo a indivíduo),
uma versão híbrida otimizada (processamento em lote) e uma versão
totalmente residente em GPU. Os experimentos, conduzidos em um conjunto
de 38 instâncias da TSPLIB com até 1002 cidades, demonstram que a
estratégia de paralelização influencia drasticamente o trade-off entre
tempo de execução e qualidade da solução. A variante híbrida otimizada
obteve os menores tempos de execução (ganho de performance mediano por
instância de 5,69× sobre a híbrida ingênua), enquanto a variante
totalmente em GPU alcançou a melhor qualidade de solução (gap médio de
0,88\%). Os resultados evidenciam que a minimização da transferência de
dados entre CPU e GPU é crítica para o desempenho, e que arquiteturas
totalmente residentes favorecem o refinamento das soluções em problemas
de otimização combinatória.

\textbf{Palavras-chave:} Problema do Caixeiro Viajante, Algoritmos
Genéticos, Computação em GPU, CUDA, Otimização Combinatória.

\clearpage

\section*{Abstract}\label{abstract}
\addcontentsline{toc}{section}{Abstract}

This work investigates the impact of different parallelization
strategies on Graphics Processing Units (GPUs) on the performance of a
memetic algorithm (Genetic Algorithm hybridized with 2-opt local search)
for the Traveling Salesman Problem (TSP). Four iso-algorithmic variants
were implemented and compared: a sequential CPU version, a naive hybrid
version (individual-by-individual transfer), an optimized hybrid version
(batch processing), and a fully GPU-resident version. Experiments
conducted on a set of 38 TSPLIB instances with up to 1002 cities
demonstrate that the parallelization strategy drastically influences the
trade-off between execution time and solution quality. The optimized
hybrid variant achieved the lowest execution times (median per-instance
Ganho de Performance of 5.69× over the naive hybrid), while the fully
GPU-resident variant achieved the best solution quality (average gap of
0.88\%). The results highlight that minimizing data transfer between CPU
and GPU is critical for performance, and that fully resident
architectures favor solution refinement in combinatorial optimization
problems.

\textbf{Keywords:} Traveling Salesman Problem, Genetic Algorithms, GPU
Computing, CUDA, Combinatorial Optimization.

\clearpage

\clearpage
\csname tableofcontents\endcsname
\clearpage

\section*{Glossário de termos}\label{glossario}
\addcontentsline{toc}{section}{Glossário de termos}

\begin{longtable}[]{@{}
  >{\raggedright\arraybackslash}p{(\linewidth - 4\tabcolsep) * \real{0.3333}}
  >{\raggedright\arraybackslash}p{(\linewidth - 4\tabcolsep) * \real{0.3333}}
  >{\raggedright\arraybackslash}p{(\linewidth - 4\tabcolsep) * \real{0.3333}}@{}}
\toprule\noalign{}
\begin{minipage}[b]{\linewidth}\raggedright
Termo em inglês
\end{minipage} & \begin{minipage}[b]{\linewidth}\raggedright
Forma adotada no texto
\end{minipage} & \begin{minipage}[b]{\linewidth}\raggedright
Observação
\end{minipage} \\
\midrule\noalign{}
\endhead
\bottomrule\noalign{}
\endlastfoot
\emph{batch} & lote & Usado para processamento de vários indivíduos em
uma única chamada. \\
\emph{benchmark} & benchmark / avaliação experimental & Mantém-se
``benchmark'' quando o foco é o protocolo comparativo. \\
\emph{crossover} & cruzamento & Preferência por ``cruzamento'', mantendo
a sigla OX quando necessário. \\
\emph{fitness} & aptidão & Quando útil, escreve-se ``aptidão
(\emph{fitness})'' na primeira ocorrência. \\
\emph{full GPU} / \emph{GPU-resident} & totalmente em GPU / residente na
GPU & Expressões usadas para a variante com dados mantidos em VRAM. \\
\emph{gap} & gap / desvio percentual em relação ao ótimo & O termo
``gap'' é preservado por ser consagrado na literatura. \\
\emph{host} & hospedeiro (CPU) & Contraposto a \emph{device} (GPU). \\
\emph{kernel launch} & lançamento de \emph{kernel} & O termo
\emph{kernel} é mantido por ser técnico e específico de CUDA. \\
\emph{offspring} & descendentes & Refere-se aos indivíduos gerados após
cruzamento e mutação. \\
\emph{overhead} & sobrecarga & Usado para custos adicionais de
comunicação, sincronização ou inicialização. \\
\emph{speedup} & ganho de performance & Preferência por esta expressão
no texto analítico. \\
\emph{swap mutation} & mutação por troca & Operador que troca a posição
de duas cidades na permutação. \\
\emph{throughput} & vazão & Preferência por ``vazão'', mantendo o inglês
quando a literatura exigir. \\
\emph{warp} & \emph{warp} (grupo de 32 \emph{threads}) & O termo inglês
é mantido, com explicação curta na primeira ocorrência. \\
\end{longtable}

\clearpage

\section*{Lista de Símbolos}\label{lista-de-simbolos}
\addcontentsline{toc}{section}{Lista de Símbolos}

\begin{longtable}[]{@{}
  >{\raggedright\arraybackslash}p{(\linewidth - 4\tabcolsep) * \real{0.3333}}
  >{\raggedright\arraybackslash}p{(\linewidth - 4\tabcolsep) * \real{0.3333}}
  >{\raggedright\arraybackslash}p{(\linewidth - 4\tabcolsep) * \real{0.3333}}@{}}
\toprule\noalign{}
\begin{minipage}[b]{\linewidth}\raggedright
Símbolo
\end{minipage} & \begin{minipage}[b]{\linewidth}\raggedright
Descrição
\end{minipage} & \begin{minipage}[b]{\linewidth}\raggedright
Ver seção
\end{minipage} \\
\midrule\noalign{}
\endhead
\bottomrule\noalign{}
\endlastfoot
\(n_{coords}\) & Número de cidades (coordenadas) da instância &
\hyperref[sec-33-selecao-instancias]{3.3} \\
\(n_{pop}\) & Tamanho da população do Algoritmo Genético &
\hyperref[sec-321-ga-base]{3.2.1} \\
\(n_{reps}\) & Número de repetições independentes por experimento &
\hyperref[sec-33-selecao-instancias]{3.3} \\
\(p_{\text{mut}}\) & Taxa de mutação &
\hyperref[sec-23-algoritmos-geneticos]{2.3} \\
\(p_c\) & Taxa de cruzamento &
\hyperref[sec-23-algoritmos-geneticos]{2.3} \\
\(k_{\text{trnmt}}\) & Tamanho do torneio na seleção &
\hyperref[sec-251-ga-2opt]{2.5.1} \\
\(n_{2\text{-opt}}\) & Número máximo de iterações da busca local 2-opt &
\hyperref[sec-251-ga-2opt]{2.5.1} \\
\(B_{\mathrm{H2D}}\) & Volume de dados transferidos Host-to-Device &
\hyperref[cap3-materiais-metodos]{Cap. 3} \\
\(B_{\mathrm{D2H}}\) & Volume de dados transferidos Device-to-Host &
\hyperref[cap3-materiais-metodos]{Cap. 3} \\
\(T\) & Tempo de execução &
\hyperref[sec-41-resultados-experimentais]{4.1} \\
\(\Delta C\) & Variação de custo em um movimento local &
\hyperref[sec-221-2opt]{2.2.1} \\
\(G = (V, E)\) & Grafo composto por vértices \(V\) e arestas \(E\) &
\hyperref[sec-21-otimizacao-tsp]{2.1} \\
\(c_{ij}\) & Custo (distância) entre as cidades \(i\) e \(j\) &
\hyperref[sec-21-otimizacao-tsp]{2.1} \\
\(f(x)\) & Função de aptidão (\emph{fitness}) de uma solução \(x\) &
\hyperref[sec-23-algoritmos-geneticos]{2.3} \\
\(P(t)\) & População de soluções na geração \(t\) &
\hyperref[sec-23-algoritmos-geneticos]{2.3} \\
\(\kappa\) & Limite de iterações para busca local &
\hyperref[sec-221-2opt]{2.2.1} e \hyperref[sec-251-ga-2opt]{2.5.1} \\
\end{longtable}

\clearpage

\listoffigures

\clearpage

\listoftables

\section{Capítulo 1 -- Introdução}\label{cap1-introducao}

\subsection{1.1 Contexto e motivação}\label{sec-11-contexto-e-motivacao}

O Problema do Caixeiro Viajante (Traveling Salesman Problem -- TSP) é um
dos problemas mais estudados em otimização combinatória
\autocite{cook2012pursuit}. Tem esse nome devido à sua história
clássica: um vendedor precisa visitar um conjunto de cidades, passando
por cada uma exatamente uma vez, e deseja minimizar a distância total
percorrida.

O objetivo é encontrar um ciclo hamiltoniano (ciclo no qual se deve
passar por todos os vértices uma vez e retornar ao vértice original).
Apesar de sua formulação simples, o TSP é NP-difícil (classe de
problemas para os quais não se conhece algoritmo polinomial eficiente)
e, na sua formulação de decisão: ``existe um tour com custo menor ou
igual a \(B\)?'', é NP-completo: não se conhece algoritmo em tempo
polinomial que o resolva em geral, embora seja fácil verificar o custo
de uma solução candidata. \textcite{cook2012pursuit} argumenta que essa
combinação de simplicidade, dificuldade teórica e rica estrutura
geométrica faz do TSP um estudo de caso central tanto para a teoria da
complexidade quanto para desenvolvimento de métodos exatos e
heurísticos.

Na prática, variantes do TSP aparecem em domínios como roteirização de
veículos, planejamento de inspeções, manufatura e testes de circuitos,
genômica e astronomia, entre outros \autocite{cook2012pursuit}. Exemplos
clássicos incluem o posicionamento e ordenamento de furos em placas de
circuito impresso, logística e problemas de mapeamento genético. Mesmo
quando modelos reais são mais complexos (com janelas de tempo, múltiplos
veículos ou restrições de capacidade), é comum validar ideias de projeto
e análise de algoritmos primeiro em instâncias clássicas do TSP,
justamente pela ampla disponibilidade de referências, testes
padronizados e de soluções ótimas conhecidas.

Paralelamente, a evolução do hardware trouxe processadores gráficos
(GPUs) como plataforma acessível para computação de alto desempenho
(HPC) \autocite{nvidia2024cuda} e com a popularização e acessibilidade,
pesquisas nestas áreas cresceram rapidamente. Modelos de Inteligência
Artificial Generativa, por exemplo, dependem fortemente desse tipo de
hardware.

GPUs oferecem milhares de núcleos relativamente simples, organizados em
um modelo de execução massivamente paralelo normalmente descrito, em
CUDA, como SIMT (\emph{Single Instruction, Multiple Threads})
\autocite{flynn1972taxonomy,nvidia2024cuda}. Isso não significa que CPUs
sejam arquiteturas ``SIMD por natureza'': CPUs modernas continuam sendo
predominantemente MIMD (\emph{Multiple Instruction, Multiple Data}) no
nível de núcleos e \emph{threads}, embora cada núcleo possa explorar
extensões vetoriais SIMD (como SSE e AVX) para acelerar operações sobre
vários dados por instrução. Em troca de um controle mais restrito nos
fluxos de processo e memória, as GPUs entregam grande largura de banda
de memória e alta taxa de operações aritméticas por segundo (FLOPs),
desde que o problema ofereça operações semelhantes que possam ser
executadas em paralelo. Isso torna GPUs particularmente atrativas para
tarefas como o cálculo de matrizes, a avaliação de grandes populações de
soluções e a aplicação de movimentos de vizinhança independentes em
heurísticas de melhoria.

Do ponto de vista das meta-heurísticas, \textcite{crainic2003parallel}
propõem uma taxonomia de paralelização que distingue estratégias
baseadas em decomposição de dados, decomposição de domínio e múltiplas
trajetórias cooperativas. A discussão detalhada dessas categorias e sua
adequação a arquiteturas de GPU é apresentada na
\hyperref[sec-261-taxonomia-paralelizacao]{Seção 2.6.1}.

O projeto desta monografia foca principalmente em algoritmos do tipo 1,
explorando paralelismo de dados em vizinhanças de \(2\)-opt e na
avaliação de populações em algoritmos genéticos, em linha com estudos
recentes sobre heurísticas para TSP em GPU
\autocite{schulz2013gpu,tsp_gpu,fujimoto2011highly}.

Ainda assim, nem todo algoritmo se beneficia automaticamente de uma
migração para GPU. Inicializações de kernels, movimentação de dados
entre CPU e GPU podem anular ganhos de paralelismo, especialmente em
problemas de porte pequeno ou em implementações que realizam pouco
trabalho por inicializações de kernel
\autocite{van2013gpu,luong2013gpu}. Limitações de memória (VRAM) devem
ser cuidadosamente calculadas para que a memória total ocupada não passe
do limite do hardware. Nesses cenários, uma comparação direta entre
versões em CPU e GPU exige, ainda, cuidados extras: deve-se isolar o
impacto da plataforma de execução sem confundir o efeito com mudanças na
lógica do algoritmo. Ao longo deste trabalho, esses desafios são
documentados explicitamente em estudos de caso, onde kernels são
chamados de forma não cautelosa e versões otimizadas (Capítulos
\hyperref[cap3-materiais-metodos]{3} e \hyperref[cap4-resultados]{4}),
bem como em notas técnicas específicas sobre \emph{overhead} de
lançamentos de kernel e padrões de redução em GPU.

Este trabalho insere-se nesse contexto, investigando como acelerar,
utilizando GPU, um Algoritmo Genético (Genetic Algorithm -- GA) híbrido
com \(2\)-opt para o TSP, de forma controlada e estatisticamente
fundamentada, usando instâncias clássicas da biblioteca TSPLIB
\autocite{reinelt1991tsplib} e um protocolo experimental com um número
adequado de repetições \(n_{reps}\) e testes estatísticos apropriados
(detalhado no \hyperref[cap3-materiais-metodos]{Capítulo 3}). A
combinação entre GA e busca local é um exemplo de algoritmo memético
amplamente estudado na literatura
\autocite{larranaga1999genetic,goldberg1989genetic}, em que operadores
evolutivos globais são complementados por heurísticas de melhoria como
\(2\)-opt \autocite{croes1958method} ou Lin--Kernighan
\autocite{lin1973efficient}. Neste Trabalho de Conclusão de Curso, será
utilizada a nomenclatura em inglês para alguns termos técnicos como
``\(2\)-opt'', ``TSPLIB'', ``GPU'', ``CPU'' e ``GA''.

\subsection{1.2 Problema de pesquisa}\label{sec-12-problema-de-pesquisa}

Muitos estudos em computação de alto desempenho comparam algoritmos em
CPU e GPU utilizando implementações que diferem não apenas na plataforma
de execução, mas também em detalhes relevantes da lógica do algoritmo
(por exemplo, operadores distintos, parâmetros diferentes ou vizinhanças
não equivalentes)
\autocite{schulz2013gpu,van2013gpu,benaini2018genetic}. Isso torna
difícil atribuir ganhos de desempenho exclusivamente ao uso da GPU, uma
vez que alterações no desenho algorítmico ou na parametrização podem,
por si só, explicar diferenças observadas. Neste trabalho, os parâmetros
e operadores são mantidos fixos entre as variantes, e as diferenças de
desempenho observadas são discutidas à luz do protocolo experimental
apresentado no \hyperref[cap3-materiais-metodos]{Capítulo 3}.

Neste trabalho, busquei implementar as variantes do algoritmo sob as
mesmas condições. A estrutura dos algoritmos é projetada para ser
idêntica e otimizada tanto para CPU quanto para GPU, de forma justa,
para todos os Algoritmos Genéticos e para o algoritmo de busca local
\(2\)-opt, variando apenas a forma como cada ``módulo'' estrutural do
algoritmo é executado (em CPU ou em GPU) e o grau de paralelismo
utilizado. A pergunta central pode ser formulada da seguinte forma:

\begin{quote}
\textbf{Como comparar, de forma fiel, o impacto de diferentes
estratégias de paralelização utilizando GPU no desempenho de um
algoritmo memético (GA+\(2\)-opt) para o Problema do Caixeiro Viajante?}
\end{quote}

Responder a essa pergunta exige, ao mesmo tempo, um desenho experimental
cuidadoso de hiperparâmetros e validade estatística (número
significativo de instâncias de problema, número adequado de repetições,
métricas de qualidade e tempo). Os parâmetros básicos do GA --- como
tamanho de população \(n_{pop}\), taxa de mutação \(p_{\text{mut}}\),
tamanho do torneio \(k_{\text{trnmt}}\), número de iterações de
\(2\)-opt por indivíduo \(n_{2\text{-opt}}\) e número de repetições por
combinação algoritmo--instância \(n_{reps}\) --- são escolhidos com base
em recomendações da literatura de algoritmos evolutivos
\autocite{eiben2015introduction,goldberg1989genetic} e em análises
específicas documentadas nos relatórios técnicos do projeto. Esses
valores e sua motivação são apresentados em detalhe no
\hyperref[cap3-materiais-metodos]{Capítulo 3}.

\subsection{1.3 Objetivo geral}\label{sec-13-objetivo-geral}

O objetivo geral deste trabalho é investigar e quantificar o impacto de
diferentes estratégias de paralelização em processadores gráficos, com
foco em:

\begin{itemize}
\tightlist
\item
  \textbf{Tempo de execução}: medir e comparar o tempo requerido para
  cada variante do algoritmo memético GA+\(2\)-opt;
\item
  \textbf{Qualidade das soluções}: avaliar o custo final das rotas
  obtidas e o \emph{gap} relativo ao ótimo conhecido, validando a
  viabilidade da solução;
\item
  \textbf{Manutenção de equivalência funcional}: garantir que as quatro
  variantes (GA-CPU, GA-Híbrido-Ingênuo, GA-Híbrido-Otimizado,
  GA-FullGPU) implementem exatamente a mesma lógica algorítmica,
  diferindo apenas na plataforma de execução e no grau de paralelismo;
\item
  \textbf{Metodologia estatística padronizada}: aplicar testes de
  hipótese apropriados (testes de normalidade, pareados e rankings)
  conforme recomendado por \textcite{demsar2006statistical};
\item
  \textbf{Alinhamento com taxonomia de paralelização}: situar as
  estratégias adotadas no contexto da taxonomia proposta por
  \textcite{crainic2003parallel} e aprofundada em
  \textcite{crainic2010parallel}.
\end{itemize}

\subsection{1.4 Objetivos
específicos}\label{sec-14-objetivos-especificos}

Para viabilizar o objetivo geral, definem-se os seguintes objetivos
específicos:

\begin{enumerate}
\def\labelenumi{\arabic{enumi}.}
\item
  \textbf{Implementar quatro variações
  ``isoalgorítmicas''}{[}\^{}isoalgorithm-disambiguation{]} de um
  Algoritmo Genético híbrido com \(2\)-opt para o TSP, que diferem
  apenas na forma e no local de execução (CPU ou GPU):
  {[}\^{}isoalgorithm-disambiguation{]}: O termo isoalgorítmico será
  utilizado neste trabalho para refletir variantes estruturalmente
  similares em sua composição.

  \begin{itemize}
  \tightlist
  \item
    uma versão executada puramente na CPU (\textbf{GA-CPU}), em que
    todas as operações (seleção, cruzamento, mutação, \(2\)-opt e
    cálculo de custo) são executadas com NumPy, servindo como base em
    CPU para as comparações. Na prática, essa versão é utilizada como
    base principalmente para instâncias de menor porte (por exemplo, com
    \(n_{coords} \leq 100\)); para instâncias maiores, a comparação de
    tempo passa a considerar como referência a versão híbrida, conforme
    discutido no Capítulo 3.
  \item
    uma versão híbrida com \(2\)-opt executada na GPU e avaliação em CPU
    (\textbf{GA-Híbrido-Ingênuo{[}\^{}ingenuity-disambiguation{]}}), que
    inicia kernels \textbf{individualmente} para cada rota e transfere
    as soluções (ou circuitos) completos entre CPU e GPU a cada
    iteração;
  \item
    uma versão híbrida otimizada (\textbf{GA-Híbrido-Otimizado}), em que
    \(2\)-opt e o cálculo de custos são executados em GPU em lote
    (\emph{batch}), reduzindo o número de lançamentos de kernel e o
    volume de dados transferidos;
  \item
    uma versão totalmente em GPU (\textbf{GA-FullGPU}), na qual
    população, operadores genéticos, busca local e avaliação permanecem
    residentes na GPU durante toda a evolução.
  \end{itemize}

  Os nomes utilizados aqui seguem as implementações
  \texttt{GeneticAlgorithmCPU}, \texttt{GeneticAlgorithmHybridNaive},
  \texttt{GeneticAlgorithmHybridOptimized} e
  \texttt{GeneticAlgorithmFullGPU} do framework experimental, cujos
  detalhes arquiteturais são apresentados no
  \hyperref[cap3-materiais-metodos]{Capítulo 3}.
  {[}\^{}ingenuity-disambiguation{]}: Na literatura, o termo ``naive''
  (ingênuo) é frequentemente utilizado para descrever implementações
  simples que não exploram otimizações avançadas.
\item
  \textbf{Selecionar um conjunto de instâncias da TSPLIB} com diferentes
  faixas de tamanho (pequenas \(n_{coords} \leq 100\), médias
  \(100 < n_{coords} \leq 400\) e grandes \(n_{coords} > 400\)). Essa
  seleção visa demonstrar os ganhos de desempenho que podem ser
  atingidos e garantir diversidade suficiente para analisar ganhos em
  escala. A divisão de faixas adotada é inspirada em estudos prévios com
  a TSPLIB \autocite{reinelt1991tsplib}, mas levemente ajustada para
  refletir o foco deste trabalho nas instâncias pequenas, médias e
  grandes selecionadas. O ajuste consiste em adotar explicitamente o
  número de coordenadas \(n_{coords} = 100\) e \(n_{coords} = 400\) como
  limites entre as faixas, em linha com o limiar prático utilizado para
  a comparação CPU/GPU e com a distribuição de tamanhos das instâncias
  efetivamente utilizadas nos experimentos.
\item
  \textbf{Definir um protocolo experimental reprodutível}, incluindo
  número de repetições por instância e algoritmo, critérios de parada,
  parâmetros do GA e limites práticos impostos pela capacidade de
  memória da GPU. O desenho desse protocolo segue recomendações da
  literatura de comparação de algoritmos e de testes estatísticos
  \autocite{demsar2006statistical} e é detalhado no
  \hyperref[cap3-materiais-metodos]{Capítulo 3}.
\item
  \textbf{Aplicar um conjunto de testes estatísticos apropriados} para
  comparação de algoritmos, incluindo testes de normalidade
  (Shapiro--Wilk), testes pareados paramétricos (teste \emph{t-pareado})
  e não paramétricos (Wilcoxon), análise de rankings em múltiplas
  instâncias (teste de Friedman com pós-teste de Nemenyi) e medidas de
  tamanho de efeito (Cohen \(d\)), conforme recomendado por
  \textcite{demsar2006statistical} e literatura correlata de teste de
  hipóteses. A metodologia estatística completa é apresentada no
  \hyperref[cap3-materiais-metodos]{Capítulo 3}.
\item
  \textbf{Analisar os resultados obtidos}, discutindo as possíveis
  condições em que cada variante algorítmica é vantajosa --- como por
  exemplo, tamanho da instância \(n_{coords}\), relação entre custos de
  comunicação e custo computacional, acesso à memória e volume de dados
  transferidos entre CPU e GPU nas direções host-to-device (H2D ou CPU
  \(\to\) GPU) e device-to-host (D2H GPU \(\to\) CPU)
  \autocite{nvidia2024cuda} ---, limites práticos do uso de GPUs para o
  contexto de recursos limitados (paralelismo em uma placa gráfica já
  defasada e limitada) e as implicações para o projeto de algoritmos de
  roteamento em cenários reais, seguindo as taxonomias de paralelização
  de meta-heurísticas propostas por \textcite{crainic2003parallel} e
  complementadas por \textcite{alba2005parallel}, além de estudos de
  caso em problemas de roteamento utilizando GPUs
  \autocite{schulz2013gpu,tsp_gpu,fujimoto2011highly}.
\end{enumerate}

\subsection{1.5 Justificativa}\label{sec-15-justificativa}

Do ponto de vista científico, o TSP continua sendo um problema de
referência para avaliar novas ideias em heurísticas e meta-heurísticas
\autocite{cook2012pursuit}. A vasta disponibilidade de estudos,
instâncias utilizadas em larga escala e muitas soluções ótimas
conhecidas, em particular as fornecidas pela TSPLIB
\autocite{reinelt1991tsplib}, na qual o projeto se baseia, permite medir
o desempenho de diferentes algoritmos tanto em termos de qualidade
quanto em tempo de execução. Ao longo deste trabalho, será utilizada a
notação \(n_{coords}\) para o número de coordenadas (cidades) de cada
instância, \(n_{pop}\) para o tamanho da população do GA, \(n_{reps}\)
para o número de repetições por combinação algoritmo--instância, \(T\)
para tempos de execução médios e \(B_{\mathrm{H2D}}\),
\(B_{\mathrm{D2H}}\) para os volumes totais de dados transferidos entre
CPU e GPU nas direções host-to-device e device-to-host, respectivamente.
Os detalhes de como esses elementos se relacionam com os limites de
memória de cada variante são discutidos no
\hyperref[cap3-materiais-metodos]{Capítulo 3}.

No campo da computação de alto desempenho, diferentes trabalhos relatam
ganhos relevantes ao portar heurísticas de melhoria e algoritmos
evolutivos para GPU, inclusive em variantes do TSP
\autocite{schulz2013gpu,tsp_gpu,fujimoto2011highly}. Esses resultados
indicam que GPUs são um hardware promissor para esse tipo de algoritmo,
mas também evidenciam um problema recorrente: em muitas comparações, as
versões em CPU e GPU diferem não apenas na plataforma de execução, mas
também em detalhes de implementação e parametrização, o que dificulta
isolar o efeito específico da paralelização.

Além disso, para que conclusões sobre desempenho sejam confiáveis, não
basta observar poucos experimentos isolados. É necessário adotar um
desenho experimental com múltiplas repetições, tanto para instâncias
quanto para, medidas de dispersão e testes de hipótese adequados, como
discutido na literatura de comparação de algoritmos
\autocite{demsar2006statistical}. Neste trabalho, esses cuidados são
incorporados ao experimento, de modo que as diferenças observadas entre
variantes em CPU e GPU possam ser atribuídas, com maior segurança, às
estratégias de paralelização avaliadas.

Este trabalho justifica-se por combinar três elementos que, em conjunto,
fortalecem a contribuição científica:

\begin{enumerate}
\def\labelenumi{\arabic{enumi}.}
\tightlist
\item
  \textbf{Implementação de heurísticas clássicas em GPU}
  (particularmente \(2\)-opt e operadores do AG), explorando o
  paralelismo de forma explícita.
\item
  \textbf{Comparação entre variantes funcionalmente equivalentes do
  algoritmo em CPU e GPU}, mantendo fixos operadores, parâmetros e
  critérios de parada, de modo a isolar o efeito da plataforma de
  execução. {[}\^{}small-variations{]} {[}\^{}small-variations{]}:
  Algumas variações dentro deste contexto são esperadas, mas de modo
  geral, o algoritmo segue a mesma estrutura. Sempre há a possibilidade
  de haver pequenas variações, mas o autor busca manter, ao máximo, a
  consistência entre as variantes.
\item
  \textbf{Aplicação de metodologia estatística rigorosa}, com múltiplas
  repetições por instância, testes de hipótese adequados e medidas de
  tamanho de efeito.
\end{enumerate}

Os capítulos seguintes discutem em que condições cada variante é
vantajosa, quais os limites práticos do uso de GPUs e recursos limitados
e como os resultados dialogam com estudos prévios de heurísticas em GPU.

\subsection{1.6 Organização do trabalho}\label{sec-16-organizacao}

Este texto está organizado da seguinte forma:

\begin{itemize}
\tightlist
\item
  \textbf{Capítulo 2 -- Fundamentação Teórica / Revisão Bibliográfica}:
  apresenta os conceitos básicos de otimização combinatória e do TSP,
  revisa heurísticas de construção e melhoria (com ênfase em \(2\)-opt),
  discute algoritmos genéticos e meta-heurísticas híbridas, introduz
  princípios de computação em GPU e paralelização de meta-heurísticas, e
  resume abordagens estatísticas para comparação de algoritmos.
\item
  \textbf{Capítulo 3 -- Materiais e Métodos}: descreve o ambiente
  computacional, a arquitetura do framework desenvolvido (incluindo
  abstrações de backend e estratégias isoalgorítmicas), os detalhes do
  Algoritmo Genético híbrido com \(2\)-opt nas quatro variantes
  consideradas, o conjunto de instâncias TSPLIB utilizado, os parâmetros
  experimentais e o protocolo estatístico adotado.
\item
  \textbf{Capítulo 4 -- Resultados}: apresenta os resultados numéricos
  dos experimentos, incluindo estatísticas por instância, análises
  agregadas por faixa de tamanho, medidas de ganho de performance e
  qualidade das soluções, além da interpretação dos testes estatísticos
  aplicados.
\item
  \textbf{Capítulo 5 -- Discussão}: interpreta criticamente os
  resultados à luz da literatura, discute as implicações dos achados
  para o projeto de algoritmos de roteamento em GPU e analisa limitações
  do estudo.
\item
  \textbf{Capítulo 6 -- Conclusões e Trabalhos Futuros}: sintetiza as
  principais contribuições do trabalho, responde explicitamente aos
  objetivos propostos e indica possíveis extensões, como a aplicação do
  framework a problemas de roteirização mais complexos.
\item
  Elementos pós-textuais, como referências, glossário, apêndices
  técnicos e anexos com tabelas completas de resultados, são
  apresentados ao final do documento, conforme normas da instituição.
\end{itemize}

\section{Capítulo 2 - Revisão
bibliográfica}\label{cap2-revisao-bibliografica}

\subsection{2. Fundamentação teórica e revisão
bibliográfica}\label{sec-2-fundamentacao-e-revisao}

Este capítulo apresenta conceitos teóricos e a bibliografia utilizada
para contextualizar o problema estudado e as escolhas metodológicas
adotadas. São discutidos o Problema do Caixeiro Viajante (TSP) no
contexto da otimização combinatória
\autocite{cook2012pursuit,rosenkrantz1977analysis}, sua importância
geral e exemplos de utilização. Também são apresentadas generalizações
do problema e sua relação com problemas de roteamento de veículos e
variantes correlatas
\autocite{raff1983routing,bodin1981state,tan2021vehicle}, heurísticas e
meta-heurísticas clássicas para o TSP
\autocite{croes1958method,lin1973efficient,voudouris1999tsp},
aprofundando-se então no tópico de algoritmos genéticos e algoritmos
meméticos \autocite{larranaga1999genetic,goldberg1989genetic},
princípios de computação em GPU
\autocite{nvidia2024cuda,schulz2013gpu,tsp_gpu}, taxonomias de
paralelização de meta-heurísticas
\autocite{crainic2003parallel,crainic2010parallel,alba2005parallel} e de
comparação estatística de algoritmos de otimização
\autocite{demsar2006statistical}.

\subsubsection{2.1 Otimização combinatória e o Problema do Caixeiro
Viajante}\label{sec-21-otimizacao-tsp}

O TSP pode ser formulado, na versão simétrica\footnote{Na versão
  assimétrica do TSP, os custos de viagem entre pares de cidades podem
  diferir dependendo da direção (ou seja, \(c_{ij} \neq c_{ji}\)). Em
  TSPs simétricos, o custo do cálculo de distâncias pode ser
  representado por uma matriz triangular que armazena apenas uma das
  metades (superior ou inferior) sem a diagonal, o que reduz o número de
  entradas únicas de \(n^2\) para \(n(n-1)/2\).}, como um problema de
encontrar um ciclo hamiltoniano de custo mínimo --- podendo este ser a
distância percorrida, combustível gasto, ou mesmo uma combinação de
fatores. Para fins deste trabalho, o custo será também chamado de
distância --- em um grafo completo não direcionado \(G = (V, E)\), no
qual cada vértice em \(V\) representa uma cidade e cada aresta
\((i, j) \in E\) possui um custo \(c_{ij} \geq 0\) entre as cidades
\(i\) e \(j\) \autocite{cook2012pursuit}. O objetivo é determinar uma
combinação das cidades que minimize a soma total dos custos, retornando
à cidade de origem. Geralmente em aplicações práticas assume-se que a
matriz de custos é métrica e euclidiana, isto é, os custos derivam de
distâncias euclidianas entre coordenadas no plano.

Do ponto de vista da complexidade, o TSP é um problema NP-difícil em sua
forma de otimização e NP-completo em sua forma de decisão
\autocite{cook2012pursuit}. Isso significa que, em geral, não se conhece
um algoritmo em tempo polinomial que resolva instâncias arbitrárias em
grande escala. Mesmo que seja relativamente fácil verificar o custo de
uma ou várias soluções candidatas, não podemos afirmar que esta é a
solução exata em tempo polinomial. Algoritmos modernos como
\texttt{Held-Karp} \autocite{held1962dynamic} com uma complexidade
temporal \(O(n^2 2^n)\) e o algoritmo \texttt{Concorde}
\autocite{cook2012pursuit} com uma complexidade temporal indefinida,
conseguem resolver instâncias com milhares de cidades, mas seu tempo de
execução cresce exponencialmente (\(O(n^2 2^n)\)), o que, embora
assintoticamente muito superior à força bruta (\(O(n!)\)), ainda o torna
impraticável para instâncias muito grandes \footnote{O algoritmo
  Held-Karp ``troca'' a complexidade temporal da força bruta \(O(n!)\)
  por uma complexidade espacial \(O(n 2^n)\), tornando sua viabilidade
  limitada a instâncias com \(n_{coords} \approx 40\) cidades. Já o
  algoritmo Concorde, baseado em técnicas de ramificação e corte, é
  capaz de resolver instâncias com até 85.900 cidades, mas seu
  desempenho depende fortemente da estrutura específica da instância
  \autocite{cook2012pursuit}.}.

Ao longo das últimas décadas, o TSP consolidou-se também como um padrão
de avaliação empírica de algoritmos, graças à disponibilidade de
coleções de instâncias padronizadas, como a TSPLIB
\autocite{reinelt1991tsplib}. Essas coleções incluem instâncias com
diferentes tamanhos, estruturas e origens (geográficas, sintéticas,
industriais), a maioria delas com soluções ótimas conhecidas e obtidas
por métodos exatos. No presente trabalho, são utilizadas algumas dessas
instâncias como base para avaliar e comparar heurísticas e variantes
paralelas de um algoritmo memético (combinações de algoritmos evolutivos
com busca local) \autocite{larranaga1999genetic,fujimoto2011highly}.

\subsubsection{2.2 Heurísticas para o TSP}\label{sec-22-heuristicas-tsp}

Devido à dificuldade de resolver instâncias grandes do TSP exatamente,
uma vasta literatura de heurísticas tem sido desenvolvida para produzir
boas soluções em tempos de computação aceitáveis
\autocite{cook2012pursuit}. De forma geral, heurísticas podem ser
agrupadas em dois grandes tipos: heurísticas de construção e heurísticas
de melhoria.

Heurísticas de construção produzem uma solução viável ``do zero'',
frequentemente seguindo regras simples, também conhecidas como
algoritmos gulosos, como escolher iterativamente o vizinho mais próximo
ou inserir cidades em posições que causem o menor aumento de custo.
Exemplos incluem a heurística do vizinho mais próximo (nearest neighbor
--- NN), heurísticas de inserção (insertion heuristics) e variantes
baseadas em árvores de extensão mínimas (Minimum Spanning Trees ---
MSTs), como o Algoritmo de Christofides. Embora rápidas e fáceis de
implementar, essas estratégias tendem a gerar soluções de qualidade
moderada, servindo sobretudo como ponto de partida para métodos mais
sofisticados. Resultados clássicos da literatura mostram que, em
instâncias métricas, heurísticas baseadas em \texttt{MST} e em
\texttt{Christofides} podem oferecer garantias de aproximação sobre o
custo ótimo de \(\frac{3}{2}B^*\), aspecto discutido em textos de
referência em logística e roteamento \autocite{simchi2005logic}.

Heurísticas de melhoria partem de uma solução inicial e aplicam
sucessivos movimentos locais que procuram reduzir seu custo. Entre
essas, as chamadas \(k\)-opt são particularmente influentes: um
movimento \(2\)-opt consiste em remover duas arestas de um circuito e
reconectar os segmentos resultantes, escolhendo a reconexão que produz
um circuito ainda viável e de menor comprimento
\autocite{croes1958method}. Em instâncias euclidianas é comum
interpretar o \(2\)-opt como um mecanismo para eliminar cruzamentos,
mas, genericamente, ele também pode melhorar circuitos que não
apresentam interseções evidentes, ao substituir pares de arestas por
combinações de menor custo \autocite{croes1958method}. Movimentos
\(3\)-opt e extensões mais complexas, como o algoritmo de
Lin--Kernighan, generalizam essa ideia \autocite{lin1973efficient}.
Estas e outras heurísticas de melhoria desempenham papel central em
muitos algoritmos modernos para TSP, tanto como procedimentos isolados
quanto como componentes de meta-heurísticas, sendo utilizadas como
procedimentos de busca local.

Neste trabalho, a heurística \(2\)-opt é utilizada como busca local
básica acoplada ao Algoritmo Genético, em linha com estudos que combinam
heurísticas de construção simples com procedimentos de melhoria mais
complexos para obter soluções de alta qualidade em tempo razoável
\autocite{voudouris1999tsp,larranaga1999genetic,fujimoto2011highly}.

\paragraph{\texorpdfstring{2.2.1 Entendendo o algoritmo
\(2\)-opt}{2.2.1 Entendendo o algoritmo 2-opt}}\label{sec-221-2opt}

A formulação básica de um movimento \(2\)-opt pode ser descrita da
seguinte forma \autocite{croes1958method}: dado um circuito
\(C = (v_0, v_1, \dots, v_{n-1}, v_0)\) e dois índices \(i\) e \(j\)
tais que \(0 \leq i < n-1\) e \(i+2 \leq j < n\), considera-se a remoção
das arestas \((v_i, v_{i+1})\) e \((v_j, v_{j+1})\), adotando-se a
convenção circular de que, quando \(j = n-1\), o sucessor de \(v_j\) é
novamente \(v_0\). Em seguida, inserem-se as arestas \((v_i, v_j)\) e
\((v_{i+1}, v_{j+1})\). Essa reconexão implica reverter a ordem do
segmento intermediário \((v_{i+1}, \dots, v_j)\) do circuito: em outras
palavras, o \(2\)-opt não apenas troca duas conexões, mas inverte o
subcaminho entre elas. Essa operação é natural no TSP simétrico
clássico, pois a rota continua viável mesmo após a reversão. Em
variantes assimétricas ou em problemas com restrições adicionais, uma
reversão ingênua pode alterar o significado das arestas ou a viabilidade
da solução, exigindo vizinhanças adaptadas. A variação de custo
associada ao movimento é dada por

\[
\Delta C = \bigl(c_{i,j} + c_{i+1,j+1}\bigr)
         - \bigl(c_{i,i+1} + c_{j,j+1}\bigr)
\]

onde \(c_{i,j}\) denota o custo (ou distância) entre as cidades \(i\) e
\(j\). Um movimento \(2\)-opt é aceitável quando \(\Delta C < 0\), isto
é, quando a substituição das arestas reduz o comprimento total do
circuito \autocite{croes1958method}.

De forma geral, a vizinhança \(2\)-opt de um circuito é formada por
todos os pares admissíveis \((i,j)\) que definem duas arestas a serem
removidas e um segmento intermediário a ser invertido. Uma iteração do
método corresponde a uma varredura dessa vizinhança. Dependendo da
variante adotada, o algoritmo pode aplicar o primeiro movimento
melhorador encontrado (\emph{first-improvement}) ou examinar toda a
passagem e aplicar apenas o melhor (\emph{best-improvement}).

Assim, quando se determina ``\(\kappa = 10\) iterações em \(2\)-opt'',
significa permitir até 10 varreduras sucessivas da vizinhança da rota
corrente. Isso não significa ``varrer 10 vezes os 10 primeiros nós'': os
índices percorrem a rota inteira em cada passagem, e a utilidade de
múltiplas iterações está no fato de que, após cada reversão, surgem
novas arestas e novas oportunidades de melhoria. Sem esse teto, o
algoritmo continuaria até atingir um ótimo local --- o que pode ser
custoso computacionalmente em circuitos grandes
\autocite{lin1973efficient}. A forma específica como esta monografia usa
essas 10 varreduras dentro do algoritmo memético é detalhada na
\hyperref[sec-25-algoritmos-memeticos]{Seção 2.5}.

De forma mais detalhada, o algoritmo \(2\)-opt clássico pode ser
descrito como um procedimento iterativo de busca local aplicada a um
único circuito. Em linhas gerais, o algoritmo segue os passos:

\begin{enumerate}
\def\labelenumi{\arabic{enumi}.}
\tightlist
\item
  Começar com um circuito viável inicial \(C\) (obtido por uma
  heurística de construção qualquer).
\item
  Definir uma ordem de varredura para pares de índices \((i,j)\) com
  \(0 \leq i < n-1\) e \(i+2 \leq j < n\). No presente trabalho,
  utiliza-se a ordem lexicográfica padrão, em que para cada \(i\),
  varia-se \(j = i+2\) até \(n-1\) (a restrição \(j \geq i+2\) evita
  movimentos triviais{[}\^{}trivial{]} e garante que o segmento a ser
  revertido tenha comprimento mínimo \autocite{croes1958method}). Quando
  \(j = n-1\), considera-se novamente o vértice inicial \(v_0\) como
  sucessor da última cidade. Cada passagem completa por essa ordem conta
  como \textbf{uma iteração} da busca local. {[}\^{}trivial{]}:
  Movimentos triviais são aqueles que não alteram efetivamente o
  circuito, como tentar reverter um segmento de \(i\) a \(i+1\), pois
  \(c_{i, i+1}=c_{i+1, i}\) é uma aresta única e sua reversão não muda o
  circuito, no caso do TSP clássico explorado neste documento. O mesmo
  não se aplica para o ATSP.
\item
  Para cada par possível \((i,j)\), calcular \(\Delta C\) conforme a
  expressão acima.
\item
  Selecionar um critério de atualização: aplicar o primeiro melhor
  movimento encontrado (\emph{first-improvement}) ou, alternativamente,
  concluir a passagem e aplicar apenas o melhor movimento da iteração
  (\emph{best-improvement}).
\item
  Repetir a varredura sobre o circuito atualizado até que nenhum
  movimento com \(\Delta C < 0\) exista ou até alcançar o limite máximo
  de iterações \(\kappa\).
\end{enumerate}

Em uma implementação ingênua, a vizinhança \(2\)-opt de um circuito com
\(n\) cidades possui ordem \(O(n^2)\) movimentos possíveis, o que
implica um custo potencialmente elevado quando todos os pares \((i,j)\)
são examinados de forma exaustiva. Diversas otimizações são discutidas
na literatura, como o uso de listas de vizinhança, estruturas de dados
para poda de movimentos claramente não promissores e estratégias de
parada antecipada \autocite{lin1973efficient,larranaga1999genetic}. No
contexto desta monografia, o foco recai principalmente na forma como
esse procedimento é paralelizado e acoplado ao Algoritmo Genético, mais
do que em otimizações finas da vizinhança em CPU.

A Figura\textasciitilde{}\ref{fig:2opt-flowchart} ilustra uma variante
clássica do algoritmo \(2\)-opt baseada em \emph{best-improvement}: o
laço externo continua enquanto houver melhorias e o número de varreduras
completas não exceder o limite \(\kappa\). Em cada iteração, todos os
pares admissíveis \((i,j)\) são examinados, o melhor movimento da
passagem é armazenado e apenas esse melhor movimento é aplicado ao final
da varredura.

\phantomsection\label{fig:2opt-flowchart}%
\begin{Shaded}
\begin{Highlighting}[]
\NormalTok{flowchart TD}
\NormalTok{    Start(["Início: Circuito C"]) {-}{-}\textgreater{} Init["varredura ← 0"]}
\NormalTok{    Init {-}{-}\textgreater{} CheckIter\{"varredura \textless{} κ?"\}}
\NormalTok{    CheckIter {-}{-}\textgreater{}|Não| End(["Fim: Retorna C"])}
\NormalTok{    CheckIter {-}{-}\textgreater{}|Sim| ResetBest["melhor\_Δ ← 0\textless{}br/\textgreater{}melhor\_i ← {-}1\textless{}br/\textgreater{}melhor\_j ← {-}1"]}
\NormalTok{    ResetBest {-}{-}\textgreater{} LoopI["Para i = 0 até n{-}2"]}
\NormalTok{    LoopI {-}{-}\textgreater{} LoopJ["Para j = i+2 até n{-}1"]}
\NormalTok{    LoopJ {-}{-}\textgreater{} CalcDelta["Calcular ΔC = (c\_ij + c\_i+1,j+1)\textless{}br/\textgreater{}{-} (c\_i,i+1 + c\_j,j+1)"]}
\NormalTok{    CalcDelta {-}{-}\textgreater{} CheckDelta\{"ΔC \textless{} melhor Δ?"\}}
\NormalTok{    CheckDelta {-}{-}\textgreater{}|Sim| StoreBest["Guardar melhor par\textless{}br/\textgreater{}(i,j) e melhor\_Δ"]}
\NormalTok{    CheckDelta {-}{-}\textgreater{}|Não| NextJ}
\NormalTok{    StoreBest {-}{-}\textgreater{} NextJ["Próximo j"]}
\NormalTok{    NextJ {-}{-}\textgreater{}|Mais j| LoopJ}
\NormalTok{    NextJ {-}{-}\textgreater{}|Fim j| NextI["Próximo i"]}
\NormalTok{    NextI {-}{-}\textgreater{}|Mais i| LoopI}
\NormalTok{    NextI {-}{-}\textgreater{}|Fim i| ApplyBest\{"melhor\_Δ \textless{} 0?"\}}
\NormalTok{    ApplyBest {-}{-}\textgreater{}|Sim| Apply["Reverter segmento\textless{}br/\textgreater{}(melhor\_i+1..melhor\_j)\textless{}br/\textgreater{}varredura ← varredura + 1"]}
\NormalTok{    ApplyBest {-}{-}\textgreater{}|Não| End}
\NormalTok{    Apply {-}{-}\textgreater{} CheckIter}
\end{Highlighting}
\end{Shaded}

\refstepcounter{figure}\label{fig:2opt-flowchart}

Figura\textasciitilde{}\ref{fig:2opt-flowchart}: Fluxograma do algoritmo
\(2\)-opt.

\begin{Shaded}
\begin{Highlighting}[]
\NormalTok{Entrada: circuito inicial C, limite κ}
\NormalTok{Saída: circuito melhorado C}

\NormalTok{varredura = 0}

\NormalTok{enquanto varredura \textless{} κ faça}
\NormalTok{    melhor\_delta = 0}
\NormalTok{    melhor\_i = {-}1}
\NormalTok{    melhor\_j = {-}1}

\NormalTok{    para i = 0 até n{-}2 faça}
\NormalTok{        para j = i+2 até n{-}1 faça}
\NormalTok{            delta = dist(i, j) + dist(i+1, j+1) {-} (dist(i, i+1) + dist(j, j+1))}

\NormalTok{            se delta \textless{} melhor\_delta então}
\NormalTok{                melhor\_delta = delta}
\NormalTok{                melhor\_i = i}
\NormalTok{                melhor\_j = j}
\NormalTok{            fim se}
\NormalTok{        fim para}
\NormalTok{    fim para}

\NormalTok{    se melhor\_delta \textless{} 0 então}
\NormalTok{        Reverter segmento C[melhor\_i+1...melhor\_j]}
\NormalTok{        varredura = varredura + 1}
\NormalTok{    senão}
\NormalTok{        interromper}
\NormalTok{    fim se}
\NormalTok{fim enquanto}

\NormalTok{retornar C}
\end{Highlighting}
\end{Shaded}

\subsubsection{2.3 Algoritmos
genéticos}\label{sec-23-algoritmos-geneticos}

Algoritmos Genéticos (Genetic Algorithms -- GAs) são meta-heurísticas
inspiradas em princípios de evolução biológica, nas quais uma população
de soluções candidatas é iterativamente modificada por operadores
análogos à seleção natural, recombinação e mutação
\autocite{goldberg1989genetic,eiben2015introduction}. Na sua forma mais
simples, um AG mantém uma população de indivíduos representando soluções
para o problema; em cada geração, indivíduos são selecionados com base
em uma função de aptidão (fitness), recombinados por operadores de
cruzamento e perturbados por operadores de mutação. A nova população
resultante substitui total ou parcialmente a anterior, e o processo se
repete até que um critério de parada seja satisfeito.

No contexto do TSP, é comum representar cada solução como uma permutação
das cidades, utilizar operadores de cruzamento especializados para rotas
--- como Order Crossover (OX), Partially Mapped Crossover (PMX), entre
outros --- e empregar mutações que preservem a viabilidade da
permutação, como trocas de posição entre duas cidades
\autocite{larranaga1999genetic}. Estudos clássicos mostram que GAs podem
produzir soluções competitivas para o TSP quando combinados com
operadores e parâmetros adequados \autocite{goldberg1989genetic}.

Do ponto de vista formal, um GA opera sobre uma população
\(P(t) = \{x_1^{(t)}, x_2^{(t)}, \dots, x_{n_{pop}}^{(t)}\}\) de
indivíduos (soluções candidatas) na geração \(t\). Cada indivíduo
\(x_i\) possui um valor de aptidão (\emph{fitness}) \(f(x_i)\) que mede
a qualidade da solução; no caso do TSP, \(f(x_i)\) corresponde ao custo
do circuito representado por \(x_i\), e busca-se minimizar esse valor. A
cada geração, o GA aplica três operadores principais
\autocite{goldberg1989genetic,eiben2015introduction}:

\begin{enumerate}
\def\labelenumi{\arabic{enumi}.}
\item
  \textbf{Seleção}: escolhe indivíduos de \(P(t)\) com probabilidade
  proporcional (ou baseada em ranking/torneio) à sua aptidão,
  favorecendo soluções de melhor qualidade. Formalmente, a probabilidade
  de selecionar \(x_i\) pode ser dada por
  \(p_i = f(x_i) / \sum_{j=1}^{n_{pop}} f(x_j)\) (seleção proporcional)
  ou por esquemas de torneio, nos quais um subconjunto aleatório de
  indivíduos compete e o melhor é escolhido.
\item
  \textbf{Cruzamento (recombinação)}: combina pares de indivíduos
  selecionados (``pais'') para produzir descendentes. No TSP, operadores
  como OX e PMX preservam a viabilidade das permutações, herdando
  subsequências de um dos pais e preenchendo as posições restantes com a
  ordem do outro \autocite{larranaga1999genetic}. A taxa de cruzamento
  \(p_c\) controla a fração de pares submetidos a esse operador.
\item
  \textbf{Mutação}: introduz pequenas perturbações aleatórias em
  indivíduos, mantendo diversidade na população. No TSP, mutações
  típicas incluem troca de duas cidades ou reversão de segmentos. A taxa
  de mutação \(p_{\text{mut}}\) controla a probabilidade de cada
  indivíduo sofrer mutação.
\end{enumerate}

Após aplicar esses operadores, forma-se a nova população \(P(t+1)\) por
meio de um esquema de substituição (geracional, com elitismo,
steady-state, etc.). A Figura\textasciitilde{}\ref{fig:ga-flowchart}
ilustra o fluxo geral de um GA
\autocite{goldberg1989genetic,eiben2015introduction}: a população evolui
iterativamente por meio de seleção (escolha de pais com base em
aptidão), cruzamento (combinação de pares de pais gerando descendentes),
e mutação (perturbações para diversidade), até que um critério de parada
seja atingido.

\begin{Shaded}
\begin{Highlighting}[]
\NormalTok{INÍCIO AG}
\NormalTok{    Gerar população inicial aleatoriamente.}
\NormalTok{    ENQUANTO critério\_de\_parada não satisfeito FAÇA}
\NormalTok{        INÍCIO}
\NormalTok{            Selecionar pais da população.}
\NormalTok{            Produzir descendentes a partir dos pais selecionados.}
\NormalTok{            Mutar os indivíduos.}
\NormalTok{            Expandir a população adicionando os descendentes.}
\NormalTok{            Reduzir a população expandida.}
\NormalTok{        FIM}
\NormalTok{    Retornar o melhor indivíduo encontrado.}
\NormalTok{FIM AG}
\end{Highlighting}
\end{Shaded}

\phantomsection\label{fig:ga-flowchart}%
\begin{Shaded}
\begin{Highlighting}[]
\NormalTok{flowchart TD}
\NormalTok{    Start(["Início"]) {-}{-}\textgreater{} InitPop["Gerar população inicial P(0)\textless{}br/\textgreater{}aleatória ou heurística"]}
\NormalTok{    InitPop {-}{-}\textgreater{} Eval0["Avaliar aptidão f(x\_i)\textless{}br/\textgreater{}para cada indivíduo"]}
\NormalTok{    Eval0 {-}{-}\textgreater{} SetT["t ← 0"]}
\NormalTok{    SetT {-}{-}\textgreater{} CheckStop\{"Critério de\textless{}br/\textgreater{}parada\textless{}br/\textgreater{}satisfeito?"\}}
\NormalTok{    CheckStop {-}{-}\textgreater{}|Sim| Output(["Fim: Retorna\textless{}br/\textgreater{}melhor solução"])}
\NormalTok{    CheckStop {-}{-}\textgreater{}|Não| Selection["Seleção: escolher pais\textless{}br/\textgreater{}de P(t) com base em f(x\_i)"]}
\NormalTok{    Selection {-}{-}\textgreater{} Crossover["Cruzamento: combinar pares\textless{}br/\textgreater{}de pais (taxa p\_c) gerando\textless{}br/\textgreater{}descendentes"]}
\NormalTok{    Crossover {-}{-}\textgreater{} Mutation["Mutação: perturbar indivíduos\textless{}br/\textgreater{}(taxa p\_mut) para diversidade"]}
\NormalTok{    Mutation {-}{-}\textgreater{} EvalNew["Avaliar aptidão dos\textless{}br/\textgreater{}novos indivíduos"]}
\NormalTok{    EvalNew {-}{-}\textgreater{} Replace["Substituição: formar P(t+1)\textless{}br/\textgreater{}combinando P(t) e descendentes\textless{}br/\textgreater{}(com ou sem elitismo)"]}
\NormalTok{    Replace {-}{-}\textgreater{} IncT["t ← t + 1"]}
\NormalTok{    IncT {-}{-}\textgreater{} CheckStop}
\end{Highlighting}
\end{Shaded}

\refstepcounter{figure}\label{fig:ga-flowchart}

Figura\textasciitilde{}\ref{fig:ga-flowchart}: Fluxograma de Algoritmo
Genético.

\begin{Shaded}
\begin{Highlighting}[]
\NormalTok{Input: Population size N, mutation rate p\_mut, crossover rate p\_cross}
\NormalTok{Output: Best solution found}

\NormalTok{Initialize population P(0) randomly}
\NormalTok{Evaluate fitness of all individuals in P(0)}
\NormalTok{t = 0}

\NormalTok{while not stop\_condition do}
\NormalTok{    P\_new = \{\}}

\NormalTok{    // Elitism: keep best individuals}
\NormalTok{    Add best of P(t) to P\_new}

\NormalTok{    while size(P\_new) \textless{} N do}
\NormalTok{        // Selection}
\NormalTok{        parent1 = TournamentSelection(P(t))}
\NormalTok{        parent2 = TournamentSelection(P(t))}

\NormalTok{        // Crossover}
\NormalTok{        if random() \textless{} p\_cross then}
\NormalTok{            child = OrderCrossover(parent1, parent2)}
\NormalTok{        else}
\NormalTok{            child = parent1}
\NormalTok{        end if}

\NormalTok{        // Mutation}
\NormalTok{        if random() \textless{} p\_mut then}
\NormalTok{            SwapMutation(child)}
\NormalTok{        end if}

\NormalTok{        Add child to P\_new}
\NormalTok{    end while}

\NormalTok{    P(t+1) = P\_new}
\NormalTok{    Evaluate fitness of P(t+1)}
\NormalTok{    t = t + 1}
\NormalTok{end while}

\NormalTok{return Best of P(t)}
\end{Highlighting}
\end{Shaded}

\subsubsection{2.4 Heurísticas
híbridas}\label{sec-24-heuristicas-hibridas}

De forma ampla, heurísticas híbridas combinam duas ou mais estratégias
de busca, procurando explorar sinergias entre métodos com forças
complementares. No contexto de problemas de roteamento e, em particular,
do TSP, é comum combinar heurísticas de construção (como vizinho mais
próximo, inserções ou heurísticas baseadas em MST) com heurísticas de
melhoria (como movimentos \(k\)-opt) ou com meta-heurísticas
populacionais (como GAs), de modo que uma componente forneça boas
soluções iniciais e outra detalhe a exploração de vizinhanças
\autocite{voudouris1999tsp}.

Diversos autores classificam heurísticas híbridas em categorias como:
(i) \textbf{hibridização em nível de solução}, na qual uma
meta-heurística utiliza sistematicamente uma busca local (ou outro
procedimento) para melhorar indivíduos candidatos; (ii)
\textbf{hibridização em nível de algoritmo}, em que diferentes técnicas
(por exemplo, um método exato e um heurístico) são combinadas em fases
distintas ou em esquemas de cooperação; e (iii) \textbf{hibridizações
multiestágio}, nas quais diferentes heurísticas atuam em momentos
distintos do processo de resolução (inicialização, intensificação,
diversificação, etc.) \autocite{larranaga1999genetic,voudouris1999tsp}.
Essas categorias ajudam a organizar o espectro de abordagens híbridas
reportadas na literatura, embora nem sempre haja consenso absoluto sobre
as fronteiras entre elas.

Uma extensão importante dessa ideia é o conceito de algoritmos
meméticos, nos quais operadores evolutivos (seleção, cruzamento,
mutação) são combinados com heurísticas de busca local aplicadas a
indivíduos da população, de forma sistemática
\autocite{larranaga1999genetic}. Em problemas de roteamento, isso
resulta em esquemas GA+LS, nos quais um GA guia a exploração global do
espaço de soluções e um procedimento de melhoria, como \(2\)-opt ou
Lin--Kernighan, refina soluções promissoras. Esses algoritmos tendem a
oferecer melhor equilíbrio entre exploração e intensificação do que GAs
``puros'', ao custo de maior tempo computacional por iteração.

Trabalhos como o de \textcite{fujimoto2011highly} exploram justamente
essa combinação GA+busca local para o TSP em ambientes de computação
paralela, motivando a adoção de uma abordagem semelhante neste estudo. A
\hyperref[sec-25-algoritmos-memeticos]{próxima seção} aprofunda o
conceito de algoritmos meméticos e a \hyperref[sec-251-ga-2opt]{subseção
2.5.1} descreve, em linhas gerais, o esquema GA+\(2\)-opt adotado nesta
monografia.

\subsubsection{2.5 Algoritmos
meméticos}\label{sec-25-algoritmos-memeticos}

Algoritmos meméticos podem ser vistos como uma família de
meta-heurísticas evolutivas que incorporam, explicitamente, operadores
de busca local ao ciclo de evolução populacional
\autocite{larranaga1999genetic}. A metáfora usual associa o termo
``meme'' a unidades de informação que se propagam e se transformam, de
modo análogo a genes, mas em um nível mais alto de organização: além da
recombinação e mutação de indivíduos, há um processo de ``aprendizado''
local que modifica soluções de forma dirigida.

Em um algoritmo memético típico, a cada geração, parte dos indivíduos
(por exemplo, os mais aptos ou uma amostra da população) é submetida a
uma busca local, que pode ser uma heurística de melhoria como \(2\)-opt,
\(3\)-opt ou Lin--Kernighan no caso do TSP. Esse procedimento permite
refinar soluções já boas, acelerando a convergência em direção a ótimos
locais de alta qualidade. O GA fornece a diversidade global por meio de
operadores de cruzamento e mutação, enquanto a busca local intensifica a
exploração em torno de regiões promissoras do espaço de soluções.

O desenho de um algoritmo memético envolve decisões importantes, como:
(i) \textbf{quais indivíduos} serão submetidos à busca local (apenas a
elite, um subconjunto aleatório, toda a população); (ii) \textbf{com que
frequência} a busca local será aplicada (toda geração, gerações
alternadas, fases específicas do processo); e (iii) \textbf{com que
intensidade} cada chamada de busca local será executada (por exemplo,
número máximo de iterações de \(2\)-opt). Essas escolhas impactam o
balanço entre qualidade das soluções e custo computacional, bem como a
diversidade mantida na população.

Estudos fundamentais, como o de Larrañaga et al., indicam que os
algoritmos meméticos são extremamente eficientes para o Problema do
Caixeiro Viajante (TSP). Estes métodos funcionam melhor quando combinam
o cruzamento de rotas com técnicas de melhoria intensiva. Neste
trabalho, utilizamos um \textbf{Algoritmo Genético adaptado para rotas},
que possui as seguintes características:

\begin{enumerate}
\def\labelenumi{\arabic{enumi}.}
\tightlist
\item
  \textbf{Representação Numérica:} Cada solução é um circuito \(C\)
  denotando uma sequência de números inteiros representando a ordem das
  cidades.
\item
  \textbf{Seleção por Torneio:} Os ``pais'' são escolhidos através de
  competições diretas entre pequenos grupos de indivíduos.
\item
  \textbf{Cruzamento Ordenado (OX):} Um operador que preserva a ordem
  relativa das cidades herdadas dos pais.
\item
  \textbf{Mutação por Troca:} Pequenas alterações feitas trocando duas
  cidades de posição na rota.
\item
  \textbf{Substituição com Elitismo:} A nova geração substitui a antiga,
  mas preserva obrigatoriamente as melhores soluções encontradas até o
  momento.
\item
  \textbf{Hibridização:} Após essas etapas, aplica-se uma rotina de
  busca local (\(2\)-opt) para refinar os resultados.
\end{enumerate}

A viabilidade de acelerar esse processo usando placas de vídeo (GPU) é
reforçada por \textcite{fujimoto2011highly}. Nesse estudo, os autores
implementaram esse algoritmo em uma GPU (modelo GTX285) utilizando
técnicas eficientes de soma paralela para gerenciar os dados. Ao
comparar o desempenho contra um processador tradicional operando com
\textbf{apenas um núcleo} (\emph{single-core}), observaram que a GPU foi
até \textbf{24,2 vezes mais rápida} na instância \texttt{ts225}. Embora
a comparação tenha sido feita contra um processador mais antigo, o
resultado mostra que problemas de tamanho médio (200-300 cidades)
conseguem tirar proveito do paralelismo massivo da GPU.

O processo formal de refinamento segue a lógica abaixo:

\begin{enumerate}
\def\labelenumi{\arabic{enumi}.}
\tightlist
\item
  Considere uma população inicial \(P(t)\).
\item
  Aplique os operadores genéticos (cruzamento e mutação) para criar
  descendentes.
\item
  Execute a função de \textbf{Busca Local Truncada}, denotada por
  \(\mathcal{L}(x, \kappa)\), sobre cada descendente selecionado para
  refinamento.
\item
  Nesta monografia, \(\mathcal{L}(x, \kappa)\) significa aplicar
  sucessivas \textbf{varreduras completas} da vizinhança \(2\)-opt da
  rota \(x\): em cada varredura, os índices \((i,j)\) percorrem toda a
  rota admissível, calcula-se \(\Delta C\) para cada par e registra-se
  apenas o melhor movimento encontrado naquela passagem.
\item
  Ao final de cada varredura, aplica-se no máximo \textbf{uma} reversão
  de segmento --- a melhor da passagem ---, desde que o melhor
  \(\Delta C\) seja negativo. Em outras palavras, a implementação da
  monografia usa a variante \emph{best-improvement} do \(2\)-opt, e é
  aqui, no contexto memético, que a escolha específica da varredura
  passa a importar metodologicamente.
\item
  O processo continua até que:

  \begin{itemize}
  \tightlist
  \item
    a rota não possa mais ser melhorada (ótimo local); \textbf{OU}
  \item
    o número máximo de varreduras (\(\kappa\)) seja atingido.
  \end{itemize}
\item
  Nos experimentos deste trabalho, usa-se
  \(\kappa = n_{2\text{-opt}} = 10\). Portanto, cada descendente recebe
  no máximo \textbf{10 varreduras completas de \(2\)-opt} sobre toda a
  sua vizinhança admissível, e não ``10 nós'' ou ``10 trocas locais''
  arbitrárias.
\item
  O parâmetro \(\kappa\) funciona como um ``freio'': ele limita o tempo
  gasto refinando cada solução para evitar que o algoritmo fique lento
  demais, garantindo um equilíbrio entre qualidade e velocidade.
\end{enumerate}

\begin{quote}
{[}!warning{]} not what I wanted \textbf{Exemplo ilustrativo.} Considere
uma rota com 6 cidades. Em uma única varredura \emph{best-improvement},
o algoritmo examina os 10 pares admissíveis \((0,2)\), \((0,3)\),
\((0,4)\), \((0,5)\), \((1,3)\), \((1,4)\), \((1,5)\), \((2,4)\),
\((2,5)\) e \((3,5)\). Suponha que o melhor movimento da primeira
passagem seja \((1,4)\); então o segmento interno dessa rota é invertido
uma vez. Na segunda varredura, os mesmos pares de índices voltam a ser
analisados, mas agora sobre \textbf{outra rota}, pois as arestas mudaram
após a reversão. É justamente por isso que faz sentido permitir até 10
varreduras: cada nova passagem reavalia uma configuração diferente do
circuito. Se, por exemplo, na terceira passagem nenhum \(\Delta C < 0\)
for encontrado, a busca local para esse indivíduo termina antes de
chegar ao limite de 10.
\end{quote}

\paragraph{\texorpdfstring{2.5.1 Algoritmo Genético +
\(2\)-opt}{2.5.1 Algoritmo Genético + 2-opt}}\label{sec-251-ga-2opt}

O esquema GA+\(2\)-opt adotado neste trabalho segue a linha de
algoritmos meméticos padrões para o TSP \autocite{larranaga1999genetic},
nos quais um Algoritmo Genético opera sobre uma população de permutações
de cidades e, em momentos específicos do ciclo evolutivo, aplica-se uma
rotina de busca local \(2\)-opt a alguns indivíduos. Em alto nível, cada
iteração (geração) do algoritmo pode ser descrita pelos seguintes
passos:

\begin{enumerate}
\def\labelenumi{\arabic{enumi}.}
\tightlist
\item
  \textbf{Inicialização (\(t=0\))}: gerar uma população inicial \(P(0)\)
  de rotas, por exemplo, a partir de permutações aleatórias ou de
  heurísticas de construção simples (vizinho mais próximo, inserções).
\item
  \textbf{Avaliação (\emph{fitness})}: calcular o custo (comprimento)
  \(f(x_i)\) de cada rota \(x_i \in P(t)\), que servirá como medida de
  aptidão.
\item
  \textbf{Seleção}: escolher indivíduos para reprodução com base em seus
  valores de aptidão, utilizando seleção por torneio de tamanho
  \(k_{\text{trnmt}}\), favorecendo soluções de menor custo.
\item
  \textbf{Cruzamento}: combinar pares de indivíduos selecionados por
  meio de Order Crossover (OX), produzindo descendentes que preservam a
  estrutura de permutação e herdam subsequências de ambos os pais.
\item
  \textbf{Mutação}: aplicar perturbações leves às rotas (trocas de
  posição de cidades, com taxa \(p_{\text{mut}}\)) para manter
  diversidade genética.
\item
  \textbf{Busca local \(2\)-opt}: para cada descendente (ou um
  subconjunto, conforme a estratégia), aplicar
  \(\mathcal{L}(x, n_{2\text{-opt}})\) --- isto é, executar até
  \(n_{2\text{-opt}}\) varreduras completas de \(2\)-opt, cada uma
  examinando toda a vizinhança admissível e aplicando apenas o melhor
  movimento encontrado --- de modo a refinar localmente as rotas
  resultantes. Esta etapa é o que transforma o GA em um algoritmo
  memético, introduzindo intensificação explícita da busca.
\item
  \textbf{Substituição}: formar a nova população \(P(t+1)\) combinando
  os melhores indivíduos de \(P(t)\) (elitismo) com os descendentes
  refinados, garantindo que as melhores soluções sejam preservadas ao
  longo das gerações.
\end{enumerate}

O processo repete-se até que um critério de parada seja satisfeito
(número máximo de gerações, estagnação da aptidão ou atingimento do
custo ótimo conhecido). A
Figura\textasciitilde{}\ref{fig:ga2opt-flowchart} ilustra o fluxo
completo do GA+\(2\)-opt memético \autocite{larranaga1999genetic}: a
busca local de \(2\)-opt (etapa destacada) é aplicada após os operadores
genéticos (seleção por torneio, cruzamento OX, mutação swap), refinando
cada descendente com até \(n_{2\text{-opt}}\) varreduras completas de
\(2\)-opt (\(\mathcal{L}(x, n_{2\text{-opt}})\)), cada uma delas
escolhendo o melhor movimento da passagem, antes da formação da nova
população. Esse acoplamento caracteriza o algoritmo como memético,
combinando exploração global (GA) e intensificação local (\(2\)-opt).

\phantomsection\label{fig:ga2opt-flowchart}%
\begin{Shaded}
\begin{Highlighting}[]
\NormalTok{flowchart TD}
\NormalTok{    Start(["Início"]) {-}{-}\textgreater{} Init["Gerar população inicial P(0)\textless{}br/\textgreater{}(qualquer tipo de heurística construtica)"]}
\NormalTok{    Init {-}{-}\textgreater{} Eval0["Avaliar f(x\_i) para cada\textless{}br/\textgreater{}indivíduo em P(0)"]}
\NormalTok{    Eval0 {-}{-}\textgreater{} SetT["gen ← 0"]}
\NormalTok{    SetT {-}{-}\textgreater{} CheckStop\{"Critério de\textless{}br/\textgreater{}parada?"\}}
\NormalTok{    CheckStop {-}{-}\textgreater{}|Sim| Output(["Fim: Retorna melhor solução"])}
\NormalTok{    CheckStop {-}{-}\textgreater{}|Não| Selection["Seleção por torneio\textless{}br/\textgreater{}(k\_trnmt) de pais"]}
\NormalTok{    Selection {-}{-}\textgreater{} Crossover["Cruzamento (OX)\textless{}br/\textgreater{}gerar descendentes"]}
\NormalTok{    Crossover {-}{-}\textgreater{} Mutation["Mutação (swap)\textless{}br/\textgreater{}taxa p\_mut"]}
\NormalTok{    Mutation {-}{-}\textgreater{} LocalSearch["Busca Local: aplicar\textless{}br/\textgreater{}n\_2opt varreduras best{-}improvement\textless{}br/\textgreater{}de 2{-}opt a cada descendente\textless{}br/\textgreater{}(ℒ(x, n\_2opt))"]}
\NormalTok{    LocalSearch {-}{-}\textgreater{} EvalNew["Avaliar f(x\_i) dos\textless{}br/\textgreater{}descendentes refinados"]}
\NormalTok{    EvalNew {-}{-}\textgreater{} Replace["Substituição: formar P(gen+1)\textless{}br/\textgreater{}com elitismo (melhores de P(gen)\textless{}br/\textgreater{}+ descendentes)"]}
\NormalTok{    Replace {-}{-}\textgreater{} IncT["gen ← gen + 1"]}
\NormalTok{    IncT {-}{-}\textgreater{} CheckStop}
\end{Highlighting}
\end{Shaded}

\refstepcounter{figure}\label{fig:ga2opt-flowchart}

Figura\textasciitilde{}\ref{fig:ga2opt-flowchart}: Fluxograma do
algoritmo GA+\(2\)-opt.

No contexto deste trabalho, essa etapa de busca local será
posteriormente mapeada para diferentes implementações em CPU e GPU
(GA-CPU, GA-Híbrido-Ingênuo, GA-Híbrido-Otimizado, GA-FullGPU), mantendo
a mesma lógica de aplicação de \(2\)-opt (mesma vizinhança, mesmos
limites de iterações \(n_{2\text{-opt}}\)), de forma a preservar o
caráter estruturalmente equivalente das variantes. Os detalhes
operacionais, incluindo a forma precisa de parametrizar
\(n_{2\text{-opt}}\), \(n_{pop}\), \(p_{\text{mut}}\) e os critérios de
parada do GA, são apresentados no
\hyperref[cap3-materiais-metodos]{Capítulo 3}.

\subsubsection{2.6 Computação em GPU e paralelização de
meta-heurísticas}\label{sec-26-gpu-metaheuristicas}

Processadores gráficos (GPUs) evoluíram, nas últimas décadas, de
dispositivos voltados principalmente para renderização gráfica para
plataformas de computação de uso geral (GPGPU - General-purpose
computing on GPUs), amplamente utilizadas em aplicações científicas e de
inteligência artificial \autocite{nvidia2024cuda}. O modelo de
programação CUDA, por exemplo, organiza o trabalho em grades
(\emph{grids}) de blocos de threads, seguindo um paradigma de execução
massivamente paralelo denominado SIMT (Single Instruction, Multiple
Threads) \autocite{nvidia2024cuda}. Ao passo que as GPUs possuem raízes
na premissa SIMD da clássica taxonomia de Flynn \autocite{flynn1972taxonomy},
a NVIDIA estruturou este novo modelo lógico para unir eficiência
vetorial à abstração de programação escalar % TODO (Orientador): Corrigido anacronismo conceitual que associava SIMT diretamente à década de 60 na taxonomia de Flynn.
. Nessa configuração, milhares de threads
executam o mesmo kernel sobre dados distintos, o que é adequado a
tarefas com alto grau de paralelismo em dados.

Internamente, as threads em uma GPU são agrupadas em \emph{warps}
(tipicamente 32 threads em arquiteturas NVIDIA), que executam uma mesma
instrução de forma coordenada segundo o modelo SIMT (\emph{Single
Instruction, Multiple Threads}) \autocite{nvidia2024cuda}. Antes de
comparar isso com CPUs, convém fixar a nomenclatura: CPUs modernas
\textbf{não} são classificadas como SIMD; arquiteturalmente, elas
continuam sendo MIMD (\emph{Multiple Instruction, Multiple Data}), pois
diferentes núcleos e \emph{threads} podem seguir fluxos independentes.
Quando se fala em ``SIMD de CPU'', fala-se apenas de extensões vetoriais
dentro de um núcleo. Já em SIMT, cada \emph{thread} possui seu próprio
estado lógico, mas o hardware agenda grupos de \emph{threads} em passo
coordenado enquanto não há divergência. Segundo a taxonomia de Flynn
\autocite{flynn1972taxonomy}, conforme introduzido na
\hyperref[sec-11-contexto-e-motivacao]{Seção 1.1}, CPUs modernas
continuam sendo predominantemente MIMD com extensões SIMD \footnote{A
  classificação arquitetural mais comum para CPUs multicore é MIMD, pois
  diferentes núcleos ou \emph{threads} podem executar fluxos de
  instruções independentes. Instruções vetoriais como SSE e AVX
  adicionam paralelismo SIMD dentro de cada núcleo, acelerando operações
  sobre vetores de dados sem alterar a classificação geral da CPU.},
enquanto GPUs maximizam o \emph{throughput} com execução massivamente
paralela orientada a dados, sendo ideais para cargas em que a mesma
computação é aplicada a grandes volumes de dados independentes.

\paragraph{2.6.1 Taxonomia de paralelização segundo Crainic e
Toulouse}\label{sec-261-taxonomia-paralelizacao}

A taxonomia proposta por \textcite{crainic2003parallel} e aprofundada em
\textcite{crainic2010parallel} classifica as estratégias de
paralelização de meta-heurísticas em três categorias principais, cada
uma com adequação distinta a arquiteturas de hardware:

\begin{enumerate}
\def\labelenumi{\arabic{enumi}.}
\tightlist
\item
  \textbf{Tipo 1: Paralelismo de Baixo Nível (Decomposição de Dados)}

  \begin{itemize}
  \tightlist
  \item
    \textbf{Descrição:} O foco é na paralelização de tarefas
    computacionalmente intensivas dentro de uma iteração da
    meta-heurística, como a avaliação da função objetivo ou a exploração
    de vizinhança. O fluxo de controle do algoritmo permanece
    centralizado.
  \item
    \textbf{Adequação:} Ideal para arquiteturas \textbf{SIMD/SIMT
    (GPUs)}. Como a mesma operação (ex: cálculo de distância) é aplicada
    a muitos dados (ex: rotas, cálculos de matrizes), o \emph{overhead}
    de inicializações de kernels e transferência é mínimo e o
    \emph{throughput} é maximizado.
  \item
    \textbf{Exemplo Real:} \textcite{tsp_gpu} utilizam paralelismo Tipo
    1 para avaliar movimentos \(2\)-opt em paralelo na GPU. Cada thread
    calcula o ganho de uma troca de arestas, permitindo explorar
    vizinhanças massivas em tempo reduzido.
  \end{itemize}
\item
  \textbf{Tipo 2: Decomposição do Domínio}

  \begin{itemize}
  \tightlist
  \item
    \textbf{Descrição:} O espaço de busca ou as variáveis de decisão são
    particionados em subconjuntos (técnica comum em VRPs no geral), e a
    meta-heurística opera em paralelo sobre essas partições.
  \item
    \textbf{Adequação:} Mais apropriado para arquiteturas \textbf{MIMD
    (Clusters ou CPUs Multicore)}. Subproblemas distintos podem ter
    tempos de convergência variados, o que causa divergência de execução
    em uma GPU (onde as 32 threads de cada warp devem ser
    sincronizadas).
  \item
    \textbf{Exemplo Real:} Em problemas de Roteamento de Veículos com
    Múltiplos Depósitos (MDVRP), diversos autores
    \autocite[@simchi2005logic]{crainic2010parallel} descrevem
    estratégias onde cada ``depósito'' é otimizado por um processo
    independente em uma CPU, com trocas ocasionais de clientes entre
    depósitos.
  \end{itemize}
\item
  \textbf{Tipo 3: Múltiplas Trajetórias (Buscas Independentes ou
  Cooperativas)}

  \begin{itemize}
  \tightlist
  \item
    \textbf{Descrição:} Múltiplas instâncias da meta-heurística (ou
    meta-heurísticas diferentes) exploram o espaço de soluções
    simultaneamente. Podem ser independentes (\emph{multi-start}) ou
    cooperativas (troca de informações).
  \item
    \textbf{Adequação:} Versátil. Funciona bem em:

    \begin{itemize}
    \tightlist
    \item
      \textbf{Clusters/CPUs} para algoritmos complexos e heterogêneos,
      através de técnicas de \texttt{multithreading} e
      \texttt{multiprocessing}, por exemplo.
    \item
      \textbf{GPUs} para muitas execuções de algoritmos leves e
      idênticos, como \emph{Island Models}. - \textbf{Combinados}:
      Processamento massivamente paralelo, conforme revisado por
      \textcite{schulz2013gpu}.

      \begin{itemize}
      \tightlist
      \item
        \textbf{Exemplo Real:} \textcite{luong2013gpu} implementam um
        modelo de ilhas em GPU, onde milhares de pequenas populações
        evoluem em paralelo (cada uma em um bloco de threads), trocando
        periodicamente seus melhores indivíduos via memória
        compartilhada.
      \end{itemize}
    \end{itemize}
  \end{itemize}
\end{enumerate}

De forma resumida, a Figura\textasciitilde{}\ref{fig:crainic-taxonomy}
coloca lado a lado esses três tipos: a partir de uma meta-heurística
sequencial, pode-se paralelizar apenas a avaliação de
vizinhança/população (Tipo 1), decompor o conjunto de variáveis em
subproblemas (Tipo 2) ou rodar várias buscas completas em paralelo,
independentes ou cooperativas (Tipo 3).

\phantomsection\label{fig:crainic-taxonomy}%
\begin{Shaded}
\begin{Highlighting}[]
\NormalTok{flowchart TB}
\NormalTok{    MH["Meta{-}heurística sequencial"] {-}{-}\textgreater{} T1["Tipo 1\textless{}br/\textgreater{}(paralelismo em dados)\textless{}br/\textgreater{}Avaliar vizinhança/população em paralelo"]}
\NormalTok{    MH {-}{-}\textgreater{} T2["Tipo 2\textless{}br/\textgreater{}(decomposição do domínio)\textless{}br/\textgreater{}Subproblemas com subconjuntos de variáveis"]}
\NormalTok{    MH {-}{-}\textgreater{} T3["Tipo 3\textless{}br/\textgreater{}(múltiplas trajetórias)\textless{}br/\textgreater{}Buscas completas em paralelo\textless{}br/\textgreater{}independentes ou cooperativas"]}
\end{Highlighting}
\end{Shaded}

\refstepcounter{figure}\label{fig:crainic-taxonomy}

Figura\textasciitilde{}\ref{fig:crainic-taxonomy}: Resumo dos três tipos
de paralelização de Crainic e Toulouse.

No contexto deste trabalho, as variantes desenvolvidas enquadram-se
predominantemente no \textbf{Tipo 1}, focando na aceleração da avaliação
de vizinhança \(2\)-opt e operadores genéticos via paralelismo de dados
massivo na GPU.

\paragraph{2.6.2 Aplicações de GPU a meta-heurísticas para o
TSP}\label{sec-262-gpu-aplicacoes-tsp}

Aplicações de GPU a meta-heurísticas para o TSP exploram principalmente
paralelismo em dados, seja na avaliação em massa de rotas em algoritmos
genéticos (por exemplo, via cálculos de redução em GPU), seja aplicando
paralelismo a movimentos de vizinhança em heurísticas de melhoria
\autocite{schulz2013gpu,tsp_gpu,fujimoto2011highly}. Outros trabalhos
investigam paralelização de Recozimento Simulado, Colônias de Formigas
(ACO) e outras meta-heurísticas em GPU, geralmente aproveitando o grande
número de threads para explorar múltiplas soluções ou vizinhanças em
cada passo da busca
\autocite{binjubier2024gpu,rey2018cpu,abdelatti2020improvedgpuheuristic}.
Esses estudos motivam o uso de GPUs como plataforma para acelerar
algoritmos meméticos, mas também evidenciam desafios relacionados à
movimentação de dados entre CPU e GPU, à escolha de granularidade
adequada de kernels, à ocupação dos multiprocessadores e às limitações
de memória (VRAM). Esses aspectos são retomados no
\hyperref[cap3-materiais-metodos]{Capítulo 3}, ao descrever o desenho
dos kernels de \(2\)-opt e as estratégias de controle de memória
adotadas neste trabalho.

\paragraph{2.6.3 Contraste arquitetural: CPU vs
GPU}\label{sec-263-cpu-vs-gpu}

As figuras \ref{fig:cpu-gpu-contrast}, \ref{fig:cpu-gpu-detail} e
\ref{fig:gpu-internal} ajudam a visualizar, de forma simples, como CPUs
e GPUs foram pensadas para resolver problemas diferentes. De maneira
geral, CPUs têm poucos núcleos mais complexos e flexíveis. Uma parte
grande da área do chip é usada para controle de fluxo, previsão de
desvios e vários níveis de cache, o que favorece programas sequenciais,
com muitas decisões e acesso irregular à memória. Nas GPUs acontece o
contrário: a maior parte da área é ocupada por muitas unidades
aritméticas simples (ALUs), organizadas em centenas ou milhares de
núcleos em paralelo. Elas abrem mão de um controle sofisticado para
ganhar vazão (\emph{throughput}) quando muitas threads executam a mesma
sequência de instruções sobre dados diferentes. Essa diferença é o
motivo pelo qual GPUs funcionam bem para tarefas com muito paralelismo
em dados, como avaliar muitas rotas de uma vez ou aplicar \(2\)-opt em
vários indivíduos em paralelo.

\begin{figure}
\centering
\includegraphics[width=0.8\linewidth,height=\textheight,keepaspectratio]{../assets/image-1.png}
\caption{Comparação simplificada entre arquitetura de CPU e GPU. Fonte:
\textcite{layerstack2023cpugpu}.}\label{fig:cpu-gpu-contrast}
\end{figure}

A Figura\textasciitilde{}\ref{fig:cpu-gpu-detail} mostra esse contraste
na divisão da área do chip. Em uma CPU típica, poucos núcleos complexos
dividem espaço com grandes caches e lógica de controle. Em uma GPU, o
desenho é o inverso: a maior parte da área é dedicada a conjuntos de
núcleos de processamento paralelos (por exemplo, \emph{CUDA cores} em
GPUs NVIDIA ou \emph{stream processors} em GPUs AMD) e à memória de alta
largura de banda. O controle é mais simples, mas o número de operações
por segundo é muito maior quando o problema é bem mapeado para esse tipo
de arquitetura \autocite{paz2011gpucpuimage}.

\begin{figure}
\centering
\includegraphics[width=0.8\linewidth,height=\textheight,keepaspectratio]{../assets/image-2.png}
\caption{Alocação de silício em CPU e GPU. Adaptado de
\textcite{paz2011gpucpuimage}.}\label{fig:cpu-gpu-detail}
\end{figure}

A Figura\textasciitilde{}\ref{fig:gpu-internal} resume como esses
recursos aparecem na organização interna de uma GPU moderna. A memória
global (DRAM) oferece grande largura de banda, mas latência alta; cada
\emph{streaming multiprocessor} (SM) possui memórias compartilhadas
menores e registradores, usados pelas threads de um mesmo bloco. As
threads são agrupadas em \emph{warps} (tipicamente 32 threads nas GPUs
NVIDIA) que executam a mesma instrução ao mesmo tempo, seguindo o modelo
SIMT \autocite{shah2023gpuarchitectureimage}. Quando as threads de um
warp seguem caminhos de controle diferentes ou acessam a memória de
forma muito irregular, parte dessa paralelização é perdida; quando o
acesso é organizado e o fluxo de controle é parecido, o ganho de
desempenho é grande.

\begin{figure}
\centering
\includegraphics[width=0.8\linewidth,height=\textheight,keepaspectratio]{../assets/image-3.png}
\caption{Organização interna de GPU moderna e modelo de programação
CUDA. Adaptado de
\textcite{shah2023gpuarchitectureimage}.}\label{fig:gpu-internal}
\end{figure}

Essas três figuras servem como pano de fundo para as decisões de projeto
dos kernels CUDA apresentados no Capítulo 3: explorar paralelismo em
dados (tipo 1) mapeando threads para indivíduos ou movimentos \(2\)-opt,
evitar ao máximo divergência dentro de um mesmo warp e organizar as
leituras e gravações de memória global para aproveitar melhor a largura
de banda disponível.

É fundamental distinguir, portanto, três níveis de paralelismo que
aparecem neste trabalho. No nível da máquina, a CPU é tratada como uma
arquitetura MIMD: diferentes núcleos podem executar tarefas distintas ou
\emph{threads} independentes. Dentro de cada núcleo, porém, a CPU pode
explorar SIMD por meio de instruções vetoriais, o que acelera trechos
específicos do cálculo de custos e de manipulação de vetores. Já a GPU,
no contexto de CUDA, é programada no modelo SIMT: o programador escreve
milhares de \emph{threads} lógicas, mas o hardware as executa em grupos
(\emph{warps}) que avançam de forma coordenada enquanto compartilham o
mesmo fluxo de instruções. Quando há divergência de controle dentro de
um \emph{warp}, os caminhos precisam ser serializados temporariamente,
reduzindo a eficiência. Essa distinção explica por que a CPU pode ser
descrita, sem contradição, como MIMD com extensões SIMD, enquanto a GPU
é mais adequadamente descrita como SIMT.

\subsubsection{2.7 Comparação estatística de algoritmos de
otimização}\label{sec-27-comparacao-estatistica}

Meta-heurísticas estocásticas, como GAs e algoritmos meméticos, produzem
resultados que variam de execução para execução devido ao uso de
aleatoriedade em diversos pontos (inicialização, seleção, mutação, entre
outros). Por esse motivo, a comparação de algoritmos de otimização não
pode se basear em uma única execução por instância, tampouco apenas em
médias simples sem análise de variabilidade. A literatura de comparação
de algoritmos enfatiza a importância de utilizar múltiplas instâncias de
teste, múltiplas repetições por combinação algoritmo--instância e
métricas que considerem tanto qualidade da solução quanto tempo de
execução \autocite{demsar2006statistical}.

\textcite{demsar2006statistical} discute procedimentos estatísticos
apropriados para comparar algoritmos em vários problemas, recomendando o
uso de testes de hipótese pareados (como o teste \emph{t-pareado} ou o
teste de Wilcoxon) quando se deseja comparar dois algoritmos em uma
coleção de instâncias, e testes baseados em rankings, como o teste de
Friedman seguido de pós-testes de Nemenyi, quando mais de dois
algoritmos são avaliados simultaneamente. Medidas de tamanho de efeito,
como o \emph{d} de Cohen, complementam a análise ao quantificar a
magnitude prática das diferenças observadas.

Neste trabalho, esses princípios gerais orientam o desenho experimental
e a análise dos resultados apresentados nos capítulos seguintes,
garantindo que as comparações entre variantes em CPU e GPU sejam
realizadas de forma estatisticamente fundamentada. De maneira mais
concreta, para cada combinação algoritmo--instância (4 algoritmos vs 38
instâncias), são executadas múltiplas repetições independentes,
registrando-se, no mínimo, medidas de qualidade da solução (custo final
ou \emph{gap} em relação ao melhor valor conhecido) e de tempo de
execução.

\paragraph{2.7.1 Comparações
pareadas}\label{sec-271-comparacoes-pareadas}

Para comparações \textbf{par a par} entre duas variantes (por exemplo,
GA-CPU versus uma versão em GPU) ao longo de várias instâncias, adota-se
uma rotina em duas etapas: (i) verificação de normalidade das
distribuições de interesse por meio do teste de Shapiro--Wilk
\autocite{shapiro1965analysis} e (ii) aplicação de um teste pareado
apropriado. Quando a hipótese de normalidade é considerada aceitável,
utiliza-se o teste \emph{t-pareado} de Student
\autocite{student1908probable}; caso contrário, recorre-se ao teste não
paramétrico de Wilcoxon para amostras pareadas
\autocite{wilcoxon1945individual}. Em ambos os casos, medidas de tamanho
de efeito, como o \emph{d} de Cohen \autocite{cohen1988statistical}, são
utilizadas para qualificar a relevância prática das diferenças
detectadas. Valores típicos de \emph{d} de Cohen são interpretados como:
\(|d| < 0.2\) (efeito negligenciável), \(0.2 \le |d| < 0.5\) (efeito
pequeno), \(0.5 \le |d| < 0.8\) (efeito médio), e \(|d| \ge 0.8\)
(efeito grande). Estas comparações seguem as diretrizes de testes de
``Statistical Comparisons of Classifiers over Multiple Data Sets''
\autocite{demsar2006statistical}, e as fórmulas são de fácil
implementação utilizando bibliotecas em \texttt{python}.

\paragraph{2.7.2 Comparações
múltiplas}\label{sec-272-comparacoes-multiplas}

Quando três ou mais algoritmos são comparados simultaneamente em um
conjunto de instâncias (como ocorre com as quatro variantes
isoalgorítmicas GA-CPU, GA-Híbrido-Ingênuo, GA-Híbrido-Otimizado e
GA-FullGPU), emprega-se o teste de Friedman \autocite{friedman1937use}
sobre os rankings médios dos algoritmos em cada problema, seguido de
pós-testes de Nemenyi \autocite{nemenyi1963distribution} quando
apropriado, conforme recomendado por \textcite{demsar2006statistical}.
Esse procedimento permite identificar, com controle de erro tipo I,
quais pares de algoritmos apresentam diferenças estatisticamente
significativas em termos de desempenho médio. O teste de Friedman
verifica a hipótese nula de que todos os algoritmos têm desempenho
equivalente, enquanto o pós-teste de Nemenyi ajusta os valores-p para
múltiplas comparações, reduzindo o risco de falsos positivos.

Os detalhes específicos de parametrização desses testes, bem como o
conjunto exato de métricas analisadas (tempo, qualidade, ganho de
performance, entre outras), são apresentados no
\hyperref[cap3-materiais-metodos]{Capítulo 3} ao descrever o protocolo
experimental, e retomados no \hyperref[cap4-resultados]{Capítulo 4} ao
discutir os resultados obtidos.

\section{Capítulo 3 -- Materiais e
Métodos}\label{cap3-materiais-metodos}

Este capítulo descreve o ambiente computacional, o \emph{framework}
experimental e o protocolo de execução utilizados para comparar as
quatro variantes do algoritmo memético GA+\(2\)-opt apresentadas na
\hyperref[sec-14-objetivos-especificos]{Seção 1.4}. O objetivo é
permitir que outros pesquisadores reproduzam, com o máximo de precisão
possível, os experimentos discutidos no
\hyperref[cap4-resultados]{Capítulo 4}, respeitando as limitações de
hardware disponíveis.

\subsection{3.1 Ambiente computacional e framework
experimental}\label{sec-31-ambiente-framework}

Esta seção detalha as características do hardware e do software
utilizados, bem como a arquitetura do framework experimental
desenvolvido para garantir a integridade e a reprodutibilidade das
comparações.

\subsubsection{3.1.1 Hardware e sistema
operacional}\label{sec-311-hardware-software}

Os experimentos foram conduzidos em um \emph{notebook} equipado com
processador \textbf{Intel Core i7-7700HQ} e uma GPU \textbf{NVIDIA
GeForce GTX 1050 Mobile}. Embora o uso deste hardware tenha sido
motivado inicialmente pela indisponibilidade de recursos mais robustos,
ele se mostrou vantajoso metodologicamente: por se tratar de uma
arquitetura de entrada (Pascal) com limitações estritas de memória e
potência, as diferenças de eficiência entre as variantes algorítmicas
tornam-se mais evidentes do que em aceleradores de alto desempenho, onde
a força bruta computacional poderia mascarar ineficiências de
implementação. Além disso, o contraste entre uma CPU relativamente capaz
(i7 quad-core) e uma GPU de entrada acentua os desafios de
\emph{offloading}, exigindo otimizações reais para obter ganhos de
performance. Ainda, há uma maior necessidade no controle de execução e
limitações. Devido ao alto custo computacional para cálculos de
algoritmos genéticos, instâncias maiores não foram testadas. Mas módulos
e \emph{safeguards} tiveram que ser implementados no teste de outras
variações algorítmicas mais rápidas como no recozimento simulado
\autocite{wu2013performance} não discutidas neste trabalho, devido ao
baixo grau de paralelismo do algoritmo para \textbf{variações do tipo 1}
\autocite{crainic2003parallel}, obtidas experimentalmente.

As especificações relevantes da GPU para este estudo são:

O desempenho computacional é estritamente limitado pelas fronteiras
físicas do dispositivo (\emph{device}). Detalha-se a seguir o impacto
matemático destas restrições no processamento:

\begin{enumerate}
\def\labelenumi{\arabic{enumi}.}
\item
  \textbf{Memória de vídeo (VRAM) e o tamanho da instância:} A unidade
  de processamento gráfico dispõe de \(4\text{GB}\) de memória global.
  Para maximizar a capacidade de leitura, exige-se que os dados sejam
  pré-alocados, o que obriga a utilização da matriz de distâncias
  completa \((n_{coords}^2)\), em detrimento de estruturas triangulares
  otimizadas \autocite{harris2007optimizing}. Ao adotar precisão dupla
  (\texttt{double} ou \texttt{float64}, 8 bytes) para mitigar erros
  numéricos, o consumo de memória da matriz de distâncias (\(M_{dist}\))
  cresce quadraticamente:
  \[Mem(M_{dist}) = n_{coords}^2 \times 8 \text{ bytes}\] Para uma
  instância de \(n_{coords}=15.112\) cidades, a matriz ocupa
  isoladamente \(\approx 1,83 \text{GB}\). Somando-se as estruturas
  auxiliares, atinge-se o teto operacional com aproximadamente
  \(n_{coords} \approx 13.000\) cidades. A utilização de precisão
  simples (\texttt{float32}) reduziria esta ocupação pela metade,
  permitindo a resolução de instâncias de maior porte.
\item
  \textbf{Gargalo de Transferência (\emph{PCIe Transfer Bottlenecks}):}
  A comunicação entre o \emph{host} (CPU) e o \emph{device} (GPU) ocorre
  através de um barramento com largura de banda de
  \(\approx 16 \text{GB/s}\), significativamente inferior à largura de
  banda interna da memória do dispositivo (\(\approx 112 \text{GB/s}\)).
  Transferências excessivas de dados penalizam o tempo total de execução
  (\(T_{total}\)), podendo anular o ganho de aceleração computacional
  (\(T_{ganho}\)), conforme a relação: \[
   T_{total} = T_{cpu} + T_{gpu} + T_{transf}
   \] Na variante \textbf{Híbrida-Otimizada}, este efeito é mitigado
  através do encadeamento de \emph{kernels}, realizando as etapas
  intensivas (busca local e cálculo de aptidão) inteiramente no
  dispositivo e retornando ao hospedeiro apenas o vetor de custos,
  minimizando o termo \(T_{transf}\).
\item
  \textbf{Ocupação e Mascaramento de Latência:} A arquitetura é composta
  por 5 Multiprocessadores de Streaming (\emph{Streaming Multiprocessors
  - SMs}), cada um com 128 cores, totalizando 640 núcleos de execução.
  Para evitar a ociosidade destes núcleos durante os ciclos de latência
  de memória (centenas de ciclos de \emph{clock}), emprega-se a técnica
  de \emph{Latency Hiding}, mantendo múltiplas \emph{threads} residentes
  prontas para execução \autocite{kirk2010programming}.

  \begin{itemize}
  \tightlist
  \item
    \textbf{Restrição Matemática:} Cada SM suporta um máximo de 2048
    \emph{threads}. Algoritmos com alta complexidade de registradores
    reduzem a capacidade de manter \emph{threads} em espera, resultando
    em subutilização dos núcleos e perda de eficiência computacional,
    desperdiçando o potencial de processamento massivo da placa.
  \end{itemize}
\end{enumerate}

\subsubsection{3.1.2 Ambiente de
software}\label{sec-312-ambiente-software}

O framework foi desenvolvido em \textbf{python 3.10}, utilizando um
conjunto de bibliotecas selecionadas para garantir a equivalência
funcional entre as implementações em CPU e GPU:

\begin{itemize}
\tightlist
\item
  \textbf{NumPy}: Utilizado para a implementação de referência em CPU
  (GA-CPU).
\item
  \textbf{CuPy}: Utilizado para as implementações aceleradas (Híbridas e
  FullGPU). A compatibilidade de API entre CuPy e NumPy foi fundamental
  para assegurar que a lógica dos algoritmos permanecesse idêntica
  (``isoalgorítmica''), alterando apenas o \emph{backend} de execução.
\item
  \textbf{SciPy e scikit-posthocs}: Empregados para a análise
  estatística dos resultados, incluindo testes de normalidade
  (Shapiro-Wilk) e testes não-paramétricos de comparação múltipla
  (Friedman e Nemenyi).
\end{itemize}

Apesar da compatibilidade da API oferecida pelo CuPy (frequentemente
referenciada como \texttt{xp} em códigos agnósticos de backend), a
implementação de operações críticas exigiu o desenvolvimento de
\textbf{kernels CUDA customizados} (\texttt{RawKernel}). Abstrações de
alto nível do CuPy, embora conveniente, introduzem \emph{overheads} de
lançamento de kernel que se tornam proibitivos quando realizados
repetidamente para cada indivíduo da população (como na busca local
2-opt). A implementação manual de kernels permite:

\begin{enumerate}
\def\labelenumi{\arabic{enumi}.}
\tightlist
\item
  \textbf{Fusão de operações}: Realizar múltiplas etapas do 2-opt em uma
  única chamada de kernel, reduzindo a latência de comunicação CPU-GPU.
\item
  \textbf{Gerenciamento fino de memória}: Otimizar o uso da memória
  compartilhada (\emph{shared memory}) e dos registradores, maximizando
  a ocupação dos \emph{SMs}.
\item
  \textbf{Paralelismo em lote}: Processar populações inteiras
  simultaneamente, explorando o paralelismo massivo da GPU. A API padrão
  do CuPy é excelente para operações ``primitivas'' de álgebra linear
  --- como multiplicar grandes matrizes de números reais ---, onde todos
  os elementos sofrem a mesma operação. Porém, o \(2\)-opt trabalha com
  uma matriz de permutações (índices inteiros) e exige que cada linha
  (indivíduo) siga um fluxo de execução próprio, em sua thread, com
  laços \texttt{while} e condições \texttt{if} independentes. Fazer isso
  via API padrão seria inviável; por isso, foi necessário desenvolver
  kernels CUDA customizados que possibilitam que cada thread da GPU
  gerencie sua própria lógica de busca local.
\end{enumerate}

Além disso, todas as operações de ponto flutuante foram padronizadas em
\textbf{precisão dupla (\texttt{float64})}. Embora GPUs de consumo (como
a GTX 1050 utilizada) tenham desempenho superior em precisão simples
(\texttt{float32}), a precisão dupla foi mandatória para garantir a
estabilidade numérica no cálculo de distâncias acumuladas em rotas
longas (instâncias com milhares de cidades), onde erros de
arredondamento poderiam levar a avaliações incorretas de melhoria
(\texttt{delta\ \textless{}\ 0}) e divergência entre as implementações
de CPU e GPU.

\subsubsection{3.1.3 Arquitetura do framework
experimental}\label{sec-313-arquitetura-framework}

Para garantir comparações justas, o código foi estruturado em uma
arquitetura modular que separa a lógica dos algoritmos da orquestração
dos experimentos. O framework opera em três camadas distintas:

\begin{enumerate}
\def\labelenumi{\arabic{enumi}.}
\tightlist
\item
  \textbf{Camada de Algoritmos:} Contém as implementações das quatro
  variantes do algoritmo memético. Todas compartilham a mesma estrutura
  de classes e os mesmos operadores genéticos, diferindo apenas na
  estratégia de gerenciamento de memória e paralelismo.
\item
  \textbf{Camada de Orquestração (Benchmark):} Responsável por carregar
  as configurações experimentais (parâmetros, instâncias), gerenciar a
  execução das repetições independentes e garantir o isolamento entre
  testes. Esta camada implementa mecanismos de \emph{checkpoint} para
  salvar o estado de cada execução, permitindo a recuperação em caso de
  falhas e a auditoria posterior dos dados.
\item
  \textbf{Camada de Análise:} Scripts dedicados ao processamento dos
  arquivos de \emph{checkpoint}, consolidação dos dados em banco de
  dados (DuckDB) e geração automática de tabelas e relatórios
  estatísticos.
\end{enumerate}

\phantomsection\label{fig:class-diagram}%
\begin{Shaded}
\begin{Highlighting}[]
\NormalTok{classDiagram}
\NormalTok{    direction TB}

\NormalTok{    class Orchestration \{}
\NormalTok{        +run\_benchmark()}
\NormalTok{        +manage\_checkpoints()}
\NormalTok{    \}}

\NormalTok{    class GeneticAlgorithmBase \{}
\NormalTok{        \textless{}\textless{}Abstract\textgreater{}\textgreater{}}
\NormalTok{        \#population\_size: int}
\NormalTok{        \#mutation\_rate: float}
\NormalTok{        \#tournament\_size: int}
\NormalTok{        \#two\_opt\_iterations: int}
\NormalTok{        +evolve()}
\NormalTok{        \#\_initialize\_population()}
\NormalTok{        \#\_select\_parents()}
\NormalTok{        \#\_create\_offspring()}
\NormalTok{        \#\_survival\_selection()}
\NormalTok{        \#\_improve\_population()*}
\NormalTok{        \#\_evaluate\_population()*}
\NormalTok{    \}}

\NormalTok{    class GeneticAlgorithmCPU \{}
\NormalTok{        \#\_improve\_population()}
\NormalTok{        \#\_evaluate\_population()}
\NormalTok{    \}}

\NormalTok{    class GeneticAlgorithmHybridNaive \{}
\NormalTok{        \#\_improve\_population()}
\NormalTok{        \#\_evaluate\_population()}
\NormalTok{    \}}

\NormalTok{    class GeneticAlgorithmHybridOptimized \{}
\NormalTok{        \#\_improve\_population()}
\NormalTok{        \#\_evaluate\_population()}
\NormalTok{    \}}

\NormalTok{    class GeneticAlgorithmFullGPU \{}
\NormalTok{        \#\_improve\_population()}
\NormalTok{        \#\_evaluate\_population()}
\NormalTok{    \}}

\NormalTok{    Orchestration {-}{-}\textgreater{} GeneticAlgorithmBase : Instantiates \& Runs}

\NormalTok{    GeneticAlgorithmBase \textless{}|{-}{-} GeneticAlgorithmCPU}
\NormalTok{    GeneticAlgorithmBase \textless{}|{-}{-} GeneticAlgorithmHybridNaive}
\NormalTok{    GeneticAlgorithmBase \textless{}|{-}{-} GeneticAlgorithmHybridOptimized}
\NormalTok{    GeneticAlgorithmBase \textless{}|{-}{-} GeneticAlgorithmFullGPU}

\NormalTok{    note for GeneticAlgorithmBase "Padrão Template Method:\textless{}br\textgreater{}Define fluxo fixo (evolve),\textless{}br\textgreater{}subclasses implementam backends\textless{}br\textgreater{}de execução específicos (*)"}
\end{Highlighting}
\end{Shaded}

\refstepcounter{figure}\label{fig:class-diagram}

Figura\textasciitilde{}\ref{fig:class-diagram}: Diagrama de classes do
framework experimental.

Esta separação assegura que as métricas de tempo e qualidade sejam
coletadas de forma consistente para todas as variantes, eliminando
vieses que poderiam surgir de implementações \emph{ad hoc} para cada
plataforma.

\subsection{3.2 Arquitetura do framework e variantes
algorítmicas}\label{sec-32-arquitetura-variantes}

A premissa central deste trabalho é a comparação
\textbf{isoalgorítmica}\footnote{A equivalência isoalgorítmica neste estudo refere-se estritamente à estrutura, hiperparâmetros globais e ordem dos operadores. Divergências na associatividade de reduções de ponto flutuante e condições de contorno de hardware impedem um determinismo perfeito entre os dispositivos.}: todas as variantes implementam estruturalmente a
mesma meta-heurística. As
diferenças residem primariamente em \emph{onde} (CPU ou GPU) e
\emph{como} (sequencial, paralelo por indivíduo ou paralelo em lote) as
operações computacionalmente intensivas são executadas. Contudo, é crucial notar que, devido às restrições de arquitetura e mitigação de \emph{timeouts} na GPU, um teto rígido (\emph{cap}) de 50 iterações foi imposto ao laço de convergência do operador de busca local \(2\)-opt nas variantes CUDA. Esta assimetria nas condições de parada rompe a equivalência determinística com a CPU, impactando diretamente no refinamento das trajetórias de busca.

\subsubsection{3.2.1 O Algoritmo Genético Base (GA +
2-opt)}\label{sec-321-ga-base}

O algoritmo base é um Algoritmo Genético (GA) geracional com elitismo,
hibridizado com uma busca local \(2\)-opt truncada. Esta combinação,
frequentemente denominada Algoritmo Memético, equilibra a exploração
global do espaço de busca (via operadores genéticos) com a exploração
local intensiva (via \(2\)-opt).

Os componentes e parâmetros fundamentais, mantidos constantes em todas
as variantes, são:

\begin{itemize}
\tightlist
\item
  \textbf{Representação:} Permutação de inteiros representando a
  sequência de cidades visitadas.
\item
  \textbf{População Inicial:} Gerada aleatoriamente. O tamanho da
  população (\(n_{pop}\)) é adaptativo, definido como
  \(2 \times n_{coords}\), onde \(n_{coords}\) é o número de cidades da
  instância.
\item
  \textbf{Seleção:} Torneio (\emph{Tournament Selection}), favorecendo
  indivíduos com menor custo (distância total).
\item
  \textbf{Cruzamento (Crossover):} \emph{Order Crossover} (OX),
  escolhido por preservar a ordem relativa das cidades e gerar
  permutações válidas.
\item
  \textbf{Mutação:} \emph{Swap Mutation}, que troca a posição de duas
  cidades aleatórias no cromossomo.
\item
  \textbf{Busca Local (2-opt):} Aplicada a cada indivíduo da população
  ao final de cada geração. Para controlar o custo computacional, a
  busca local é \textbf{truncada}: limita-se a um número fixo de
  iterações de melhoria, entendido aqui como \textbf{10 varreduras
  completas da vizinhança por indivíduo}, em vez de buscar o ótimo local
  completo (2-opt).
\item
  \textbf{Critérios de Parada:} O algoritmo encerra sua execução se
  atingir o ótimo conhecido (com tolerância de 1\%), se não houver
  melhoria na melhor solução por um número de gerações definido pela
  ``paciência'' (\(2 \times \sqrt{n_{coords}}\)), ou se atingir o limite
  máximo de gerações (\(2 \times n_{coords} \times \sqrt{n_{coords}}\)).
\end{itemize}

\subsubsection{3.2.2 Estratégias de Paralelização (As 4
Variantes)}\label{sec-322-variantes}

Para investigar o impacto da GPU, o algoritmo base foi instanciado em
quatro variantes distintas, representando diferentes níveis de
utilização do hardware:

\begin{enumerate}
\def\labelenumi{\arabic{enumi}.}
\item
  \textbf{GeneticAlgorithmCPU (Baseline):} Execução inteiramente na CPU
  utilizando NumPy. No nível arquitetural, a CPU permanece sendo a
  plataforma MIMD de referência; o que ocorre aqui é o aproveitamento
  pontual de vetorização SIMD dentro dos núcleos para acelerar partes da
  avaliação dos movimentos 2-opt. Já o processamento dos indivíduos, no
  nível do algoritmo, ocorre essencialmente de forma sequencial. Esta
  variante serve como linha de base para medir o ganho de performance
  absoluto.
\item
  \textbf{GeneticAlgorithmHybridNaive (Híbrido Ingênuo):} Mantém a
  população e os operadores genéticos na CPU, mas transfere cada
  indivíduo para a GPU para executar a busca local 2-opt. Esta abordagem
  é considerada ``ingênua'' pois realiza transferências de memória
  (Host-to-Device e Device-to-Host) para \emph{cada indivíduo} em
  \emph{cada geração}, expondo o gargalo da largura de banda do
  barramento PCIe.
\item
  \textbf{GeneticAlgorithmHybridOptimized (Híbrido Otimizado):} Também
  mantém a lógica genética na CPU, mas otimiza a comunicação enviando a
  população inteira (ou grandes lotes) para a GPU de uma única vez. O
  kernel 2-opt processa todos os indivíduos em paralelo na GPU, e os
  resultados são retornados em lote. Esta estratégia visa amortizar o
  custo de latência das transferências de memória.
\item
  \textbf{GeneticAlgorithmFullGPU (Totalmente em GPU):} Implementa a
  estratégia de ``GPU residente''. A população é inicializada na VRAM e
  permanece lá durante todo o processo evolutivo. Todos os operadores
  (seleção, cruzamento, mutação e 2-opt) são executados via kernels CUDA
  (através do CuPy), eliminando quase completamente a necessidade de
  comunicação com a CPU, exceto para coleta de estatísticas e
  verificação de critérios de parada.
\end{enumerate}

\subsection{3.3 Seleção de instâncias e protocolo
experimental}\label{sec-33-selecao-instancias}

Para garantir a relevância estatística dos resultados, foram
selecionadas 38 instâncias da biblioteca \textbf{TSPLIB}
\autocite{reinelt1991tsplib}, variando de 51 a 1002 cidades. O conjunto
inclui problemas com diferentes características geométricas, permitindo
avaliar a robustez dos algoritmos em diversos cenários. As instâncias
selecionadas estão listadas na
Tabela\textasciitilde{}\ref{tbl:selected-instances}, categorizadas por
faixa de tamanho.

\begin{longtable}[]{@{}
  >{\raggedright\arraybackslash}p{(\linewidth - 4\tabcolsep) * \real{0.3077}}
  >{\centering\arraybackslash}p{(\linewidth - 4\tabcolsep) * \real{0.3846}}
  >{\raggedright\arraybackslash}p{(\linewidth - 4\tabcolsep) * \real{0.3077}}@{}}
\caption{Instâncias da TSPLIB selecionadas para o estudo
\label{tbl:selected-instances}}\tabularnewline
\toprule\noalign{}
\begin{minipage}[b]{\linewidth}\raggedright
Categoria
\end{minipage} & \begin{minipage}[b]{\linewidth}\centering
Faixa de Tamanho (\(n_{coords}\))
\end{minipage} & \begin{minipage}[b]{\linewidth}\raggedright
Instâncias Selecionadas
\end{minipage} \\
\midrule\noalign{}
\endfirsthead
\toprule\noalign{}
\begin{minipage}[b]{\linewidth}\raggedright
Categoria
\end{minipage} & \begin{minipage}[b]{\linewidth}\centering
Faixa de Tamanho (\(n_{coords}\))
\end{minipage} & \begin{minipage}[b]{\linewidth}\raggedright
Instâncias Selecionadas
\end{minipage} \\
\midrule\noalign{}
\endhead
\bottomrule\noalign{}
\endlastfoot
\textbf{Pequenas} & \(n_{coords} \le 100\) & eil51, berlin52, st70,
eil76, pr76, rat99, kroA100, kroB100, kroC100, kroD100, kroE100,
rd100 \\
\textbf{Médias} & \(100 < n_{coords} \le 400\) & eil101, lin105, pr107,
pr124, bier127, ch130, pr136, pr144, ch150, kroA150, kroB150, pr152,
u159, rat195, d198, kroA200, ts225, pr264, pr299, lin318, rd400 \\
\textbf{Grandes} & \(n_{coords} > 400\) & fl417, pr439, pcb442, rat783,
pr1002 \\
\end{longtable}

O protocolo de execução seguiu as diretrizes para comparação de
algoritmos estocásticos:

\begin{itemize}
\tightlist
\item
  \textbf{Repetições:} Cada par (algoritmo, instância) foi executado
  \textbf{26 vezes} independentemente.
\item
  \textbf{Sementes Aleatórias:} As sementes foram fixadas e registradas
  para garantir a reprodutibilidade, mas variaram entre as 26 repetições
  para amostrar adequadamente o comportamento estocástico.
\item
  \textbf{Limitação da CPU:} Devido ao tempo proibitivo de execução, a
  variante \textbf{GeneticAlgorithmCPU} foi executada apenas para
  instâncias com \(N \le 100\). Para instâncias maiores, as comparações
  de ganho de performance tomam como base a variante híbrida mais
  simples ou comparam as variantes GPU entre si.
\end{itemize}

\subsection{3.4 Protocolo de análise
estatística}\label{sec-34-protocolo-estatistico}

A análise dos resultados adota a metodologia recomendada por
\textcite{demsar2006statistical} para comparação de múltiplos
classificadores/algoritmos em múltiplos datasets:

\begin{enumerate}
\def\labelenumi{\arabic{enumi}.}
\tightlist
\item
  \textbf{Verificação de Normalidade:} Aplicação do teste de
  \textbf{Shapiro-Wilk} nas distribuições de tempo e qualidade das
  soluções.
\item
  \textbf{Testes Pareados:} Para comparações diretas entre duas
  variantes (ex: HybridNaive vs.~FullGPU), utiliza-se o \textbf{Teste t
  pareado} (se a distribuição for normal) ou o teste de \textbf{Wilcoxon
  Signed-Rank} (se não normal).
\item
  \textbf{Comparação Múltipla:} Para comparar as quatro variantes
  simultaneamente, aplica-se o teste de \textbf{Friedman}
  (não-paramétrico) para detectar se há diferenças significativas nos
  rankings dos algoritmos.
\item
  \textbf{Análise Post-Hoc:} Em caso de rejeição da hipótese nula no
  teste de Friedman, utiliza-se o teste de \textbf{Nemenyi} para
  identificar quais pares de algoritmos diferem significativamente entre
  si.
\item
  \textbf{Tamanho de Efeito:} O \textbf{d de Cohen} é calculado para
  quantificar a magnitude da diferença de desempenho, permitindo
  distinguir entre melhorias estatisticamente significativas mas
  irrelevantes na prática, e melhorias com impacto real no tempo de
  execução.
\end{enumerate}

\phantomsection\label{fig:stat-flowchart}%
\begin{Shaded}
\begin{Highlighting}[]
\NormalTok{flowchart TD}
\NormalTok{    Start[("Início: Escolher Teste Estatístico")] {-}{-}\textgreater{} DataType\{Tipo de Dado?\}}

\NormalTok{    DataType {-}{-}\textgreater{}|Contínuo| SampleSize\{Tamanho da Amostra?\}}
\NormalTok{    DataType {-}{-}\textgreater{}|Ordinal/Rankeado| NonParam[Testes Não{-}Paramétricos]:::nonparametric}
\NormalTok{    DataType {-}{-}\textgreater{}|Categórico| ChiSquare[Testes Qui{-}Quadrado]:::other}

\NormalTok{    SampleSize {-}{-}\textgreater{}|n \textless{} 20\textless{}br/\textgreater{}Amostra Pequena| SmallN[Obrigatório\textless{}br/\textgreater{}Teste de Normalidade]:::warning}
\NormalTok{    SampleSize {-}{-}\textgreater{}|20 \textless{}= n \textless{} 30\textless{}br/\textgreater{}Zona Cinzenta| MediumN[Recomendado\textless{}br/\textgreater{}Teste de Normalidade]:::decision}
\NormalTok{    SampleSize {-}{-}\textgreater{}|n \textgreater{}= 30\textless{}br/\textgreater{}Amostra Grande| LargeN[TCL Aplica{-}se\textless{}br/\textgreater{}Teste Recomendado]:::decision}

\NormalTok{    SmallN {-}{-}\textgreater{} Normality\{Shapiro{-}Wilk\textless{}br/\textgreater{}p \textgreater{}= 0.05?\}:::decision}
\NormalTok{    MediumN {-}{-}\textgreater{} Normality}
\NormalTok{    LargeN {-}{-}\textgreater{} Normality}

\NormalTok{    Normality {-}{-}\textgreater{}|Normal| Parametric[Testes Paramétricos]:::parametric}
\NormalTok{    Normality {-}{-}\textgreater{}|Não{-}Normal| NonParam}

\NormalTok{    Parametric {-}{-}\textgreater{} Groups\{Número de Grupos?\}:::decision}
\NormalTok{    Groups {-}{-}\textgreater{}|1 grupo| OneSample[Teste t de uma amostra]:::parametric}
\NormalTok{    Groups {-}{-}\textgreater{}|2 grupos| TwoSample[Pareado ou Independente?]:::decision}
\NormalTok{    Groups {-}{-}\textgreater{}|\textgreater{}=3 grupos| ANOVA[ANOVA / Kruskal{-}Wallis]:::parametric}

\NormalTok{    NonParam {-}{-}\textgreater{} GroupsNP\{Número de Grupos?\}:::decision}
\NormalTok{    GroupsNP {-}{-}\textgreater{}|2 pareados| Wilcoxon[Wilcoxon Signed{-}Rank]:::nonparametric}
\NormalTok{    GroupsNP {-}{-}\textgreater{}|2 independentes| MannWhitney[Mann{-}Whitney U]:::nonparametric}
\NormalTok{    GroupsNP {-}{-}\textgreater{}|\textgreater{}=3 grupos| Friedman[Friedman / Kruskal{-}Wallis]:::nonparametric}
\end{Highlighting}
\end{Shaded}

\refstepcounter{figure}\label{fig:stat-flowchart}

Figura\textasciitilde{}\ref{fig:stat-flowchart}: Fluxograma de decisão
para testes estatísticos.

\phantomsection\label{fig:results-flowchart}%
\begin{Shaded}
\begin{Highlighting}[]
\NormalTok{flowchart TD}
\NormalTok{    \%\% Nodes}
\NormalTok{    Start([Início: Análise de Resultados])}

\NormalTok{    subgraph S1 ["Estrato 1: Pequenos (n\_coords \textless{} 100)"]}
\NormalTok{        S1\_Test[Friedman Test]}
\NormalTok{        S1\_Result\{p = 0.086\}}
\NormalTok{        S1\_Conc[Inconclusivo\textless{}br/\textgreater{}Poder Estatístico Baixo]}
\NormalTok{    end}

\NormalTok{    subgraph S2 ["Estrato 2: Todos GPU{-}Only"]}
\NormalTok{        S2\_Test[Friedman Test]}
\NormalTok{        S2\_Result\{p \textless{} 0.001\}}
\NormalTok{        S2\_Nemenyi[Nemenyi Post{-}Hoc]}
\NormalTok{        S2\_Conc["FullGPU \textgreater{} Híbridos\textless{}br/\textgreater{}(Diferença Significativa)"]}
\NormalTok{    end}

\NormalTok{    subgraph Perf [Performance \& Speedup]}
\NormalTok{        Speedup\_HO[HybridOptimized: \textasciitilde{}5.69×]}
\NormalTok{        Speedup\_FG[FullGPU: \textasciitilde{}1.2x]}
\NormalTok{        Tradeoff[Trade{-}off: HybridOptimized = Velocidade\textless{}br/\textgreater{}FullGPU = Qualidade]}
\NormalTok{    end}

\NormalTok{    \%\% Flow}
\NormalTok{    Start {-}{-}\textgreater{} S1}
\NormalTok{    S1\_Test {-}{-}\textgreater{} S1\_Result {-}{-}\textgreater{}|Não Sig.| S1\_Conc}

\NormalTok{    Start {-}{-}\textgreater{} S2}
\NormalTok{    S2\_Test {-}{-}\textgreater{} S2\_Result {-}{-}\textgreater{}|Sig.| S2\_Nemenyi {-}{-}\textgreater{} S2\_Conc}

\NormalTok{    S2\_Conc {-}.{-}\textgreater{} Tradeoff}
\NormalTok{    Speedup\_HO {-}{-}\textgreater{} Tradeoff}
\NormalTok{    Speedup\_FG {-}{-}\textgreater{} Tradeoff}
\end{Highlighting}
\end{Shaded}

\refstepcounter{figure}\label{fig:results-flowchart}

Figura\textasciitilde{}\ref{fig:results-flowchart}: Fluxo de análise dos
resultados.

\section{Capítulo 4 -- Resultados}\label{cap4-resultados}

Este capítulo apresenta os resultados experimentais obtidos a partir da
execução das quatro variantes do algoritmo memético GA+\(2\)-opt (CPU,
Híbrido-Ingênuo, Híbrido-Otimizado e FullGPU) sobre o conjunto de 38
instâncias da TSPLIB. A análise é dividida em três partes principais:

\begin{enumerate}
\def\labelenumi{\arabic{enumi}.}
\tightlist
\item
  uma avaliação detalhada por instância, destacando o comportamento para
  diferentes tamanhos de problema;
\item
  uma análise agregada, focada nas métricas de ganho de performance e
  qualidade média das soluções; e
\item
  a aplicação dos testes estatísticos descritos no Capítulo 3 para
  validar as diferenças de desempenho observadas.
\end{enumerate}

\subsection{4.1 Resultados
Experimentais}\label{sec-41-resultados-experimentais}

Os experimentos foram conduzidos conforme o protocolo definido na Seção
3.3, com 26 repetições independentes para cada par algoritmo-instância.
Para as instâncias pequenas (\(n_{coords} \le 100\)), todas as quatro
variantes foram executadas. Para instâncias médias e grandes
(\(n_{coords} > 100\)), a variante puramente em CPU foi omitida devido
ao tempo de execução proibitivo, mantendo-se a comparação entre as três
variantes aceleradas por GPU.

As métricas primárias coletadas foram:

\begin{itemize}
\tightlist
\item
  \textbf{Tempo Médio de Execução (s):} Tempo total para completar a
  execução por repetição.
\item
  \textbf{Gap Médio (\%):} Desvio percentual do custo da melhor solução
  encontrada em relação ao ótimo conhecido da instância.
\item
  \textbf{Ganho de Performance:} Razão entre o tempo de execução da
  variante de referência (CPU ou Híbrido-Ingênuo) e a variante avaliada.
\end{itemize}

\subsection{4.2 Análise por
Instância}\label{sec-42-analise-por-instancia}

A Tabela\textasciitilde{}\ref{tbl:results-detailed} apresenta os
resultados detalhados para uma seleção representativa de instâncias,
cobrindo as três faixas de tamanho definidas.

\begin{longtable}[]{@{}
  >{\raggedright\arraybackslash}p{(\linewidth - 12\tabcolsep) * \real{0.1212}}
  >{\centering\arraybackslash}p{(\linewidth - 12\tabcolsep) * \real{0.1515}}
  >{\raggedright\arraybackslash}p{(\linewidth - 12\tabcolsep) * \real{0.1212}}
  >{\centering\arraybackslash}p{(\linewidth - 12\tabcolsep) * \real{0.1515}}
  >{\centering\arraybackslash}p{(\linewidth - 12\tabcolsep) * \real{0.1515}}
  >{\centering\arraybackslash}p{(\linewidth - 12\tabcolsep) * \real{0.1515}}
  >{\centering\arraybackslash}p{(\linewidth - 12\tabcolsep) * \real{0.1515}}@{}}
\caption{Resultados detalhados por instância (Seleção)
\label{tbl:results-detailed}}\tabularnewline
\toprule\noalign{}
\begin{minipage}[b]{\linewidth}\raggedright
Instância
\end{minipage} & \begin{minipage}[b]{\linewidth}\centering
Tamanho (\(n_{coords}\))
\end{minipage} & \begin{minipage}[b]{\linewidth}\raggedright
Algoritmo
\end{minipage} & \begin{minipage}[b]{\linewidth}\centering
Custo Médio (\(\pm \sigma\))
\end{minipage} & \begin{minipage}[b]{\linewidth}\centering
Gap Médio (\%)
\end{minipage} & \begin{minipage}[b]{\linewidth}\centering
Tempo Médio (s)
\end{minipage} & \begin{minipage}[b]{\linewidth}\centering
Ganho de Perf.
\end{minipage} \\
\midrule\noalign{}
\endfirsthead
\toprule\noalign{}
\begin{minipage}[b]{\linewidth}\raggedright
Instância
\end{minipage} & \begin{minipage}[b]{\linewidth}\centering
Tamanho (\(n_{coords}\))
\end{minipage} & \begin{minipage}[b]{\linewidth}\raggedright
Algoritmo
\end{minipage} & \begin{minipage}[b]{\linewidth}\centering
Custo Médio (\(\pm \sigma\))
\end{minipage} & \begin{minipage}[b]{\linewidth}\centering
Gap Médio (\%)
\end{minipage} & \begin{minipage}[b]{\linewidth}\centering
Tempo Médio (s)
\end{minipage} & \begin{minipage}[b]{\linewidth}\centering
Ganho de Perf.
\end{minipage} \\
\midrule\noalign{}
\endhead
\bottomrule\noalign{}
\endlastfoot
\textbf{berlin52} & 52 & CPU & $7554,65 \pm 22,93$ & 0,17 & 25,82 & 1,00x \\
& & HybridNaive & $7547,31 \pm 14,45$ & 0,07 & 0,57 & 45,3x \\
& & HybridOptimized & $7557,85 \pm 23,91$ & 0,21 & 0,11 & 234,7x \\
& & FullGPU & $7542,00 \pm 0,00$ & 0,00 & 0,25 & 103,3x \\
\textbf{kroA100} & 100 & CPU & $21382,88 \pm 62,04$ & 0,47 & 185,08 &
1,00x \\
& & HybridNaive & $21382,38 \pm 63,30$ & 0,47 & 1,81 & 102,3x \\
& & HybridOptimized & $21364,23 \pm 65,23$ & 0,39 & 0,30 & 616,9x \\
& & FullGPU & $21285,23 \pm 4,90$ & 0,02 & 1,37 & 135,1x \\
\textbf{pr1002} & 1002 & HybridNaive & $279875,50 \pm 3210,15$ & 8,04 &
696,05 & 1,00x \\
& & HybridOptimized & $279184,62 \pm 2789,63$ & 7,77 & 374,50 & 1,86x \\
& & FullGPU & $269776,77 \pm 743,93$ & \textbf{4,14} & 758,40 & 0,92x \\
\end{longtable}

Para instâncias pequenas como \texttt{berlin52} e \texttt{kroA100},
observa-se um ganho de performance massivo da variante
\textbf{HybridOptimized} em relação à CPU (chegando a mais de 600x no
caso de \texttt{kroA100}). Isso demonstra a eficiência da paralelização
em lote do \(2\)-opt, que amortiza o custo de lançamento de kernels. A
variante \textbf{FullGPU}, embora muito mais rápida que a CPU, apresenta
um \emph{overhead} inicial maior que a HybridOptimized para problemas
muito pequenos, devido à gestão de memória residente.

Para instâncias grandes como \texttt{pr1002}, o cenário muda. A variante
\textbf{HybridOptimized} mantém a liderança em tempo de execução (1.86x
mais rápida que a Naive), mas a variante \textbf{FullGPU} destaca-se na
\textbf{qualidade da solução}, obtendo um gap médio de 4,14\% contra
7,77\% da otimizada. Curiosamente, o tempo da FullGPU em instâncias
muito grandes aproxima-se ou até excede o da variante Naive em alguns
casos. Na ausência de dados empíricos extraídos por ferramentas de \emph{profiling}
de \emph{hardware}, as causas exatas dessa degradação de velocidade em instâncias
grandes permanecem em aberto. Os fatores podem estar relacionados tanto à
natureza do algoritmo genético escolhido quanto a limitações arquiteturais da
GPU, e a investigação detalhada dessas hipóteses constitui trabalho futuro
(ver Seção~\ref{sec-61-trabalhos-futuros}). Ainda assim, a permanência dos
dados no dispositivo evita as barreiras de latência do barramento PCIe,
favorecendo uma exploração ininterrupta que converge estatisticamente para
melhores ótimos locais.

\subsection{4.3 Análise Agregada e
Trade-offs}\label{sec-43-analise-agregada}

Ao agregar os resultados de todas as 38 instâncias, emergem padrões
claros de desempenho que caracterizam o \emph{trade-off} entre
velocidade e qualidade.

\begin{longtable}[]{@{}
  >{\raggedright\arraybackslash}p{(\linewidth - 6\tabcolsep) * \real{0.2105}}
  >{\centering\arraybackslash}p{(\linewidth - 6\tabcolsep) * \real{0.2632}}
  >{\centering\arraybackslash}p{(\linewidth - 6\tabcolsep) * \real{0.2632}}
  >{\centering\arraybackslash}p{(\linewidth - 6\tabcolsep) * \real{0.2632}}@{}}
\caption{Médias Agregadas (Todas as Instâncias)
\label{tbl:results-aggregated}}\tabularnewline
\toprule\noalign{}
\begin{minipage}[b]{\linewidth}\raggedright
Algoritmo
\end{minipage} & \begin{minipage}[b]{\linewidth}\centering
Tempo Médio (s)
\end{minipage} & \begin{minipage}[b]{\linewidth}\centering
Gap Médio ($\pm \sigma$\%)
\end{minipage} & \begin{minipage}[b]{\linewidth}\centering
Ganho de Perf. Mediano
\end{minipage} \\
\midrule\noalign{}
\endfirsthead
\toprule\noalign{}
\begin{minipage}[b]{\linewidth}\raggedright
Algoritmo
\end{minipage} & \begin{minipage}[b]{\linewidth}\centering
Tempo Médio (s)
\end{minipage} & \begin{minipage}[b]{\linewidth}\centering
Gap Médio ($\pm \sigma$\%)
\end{minipage} & \begin{minipage}[b]{\linewidth}\centering
Ganho de Perf. Mediano
\end{minipage} \\
\midrule\noalign{}
\endhead
\bottomrule\noalign{}
\endlastfoot
\textbf{HybridNaive} & 47,08 & $1,28 \pm 1,68$ & 1,00x \\
\textbf{HybridOptimized} & 18,94 & $1,30 \pm 1,78$ &
\textbf{5,69x} \\
\textbf{FullGPU} & 43,51 & $\mathbf{0,88 \pm 1,14}$ & 1,21x \\
\end{longtable}

% TODO (Ação 1.1): TABELAS INVALIDADAS — Os dados acima foram coletados com aritmética
% float32 (imprecisa) e assimetria paramétrica (FullGPU: 50 iterações de 2-opt vs 10
% para as demais). Os valores de gap e σ devem ser recalculados após a conversão para
% int32 (Ação 3.1, concluída) e a normalização do parâmetro two_opt_iterations.
% Re-execução pendente.

A análise da Tabela\textasciitilde{}\ref{tbl:results-aggregated} revela
duas conclusões fundamentais:

\begin{enumerate}
\def\labelenumi{\arabic{enumi}.}
\tightlist
\item
  \textbf{HybridOptimized é a campeã de velocidade:} Com um ganho de
  performance mediano por instância de \textbf{5,69x} sobre a versão
  ingênua, esta variante comprova que a estratégia de processamento em
  lote (\emph{batch processing}) é a mais eficaz para reduzir o tempo
  total de execução. Ao enviar a população inteira para a GPU de uma só
  vez, minimiza-se a latência de comunicação PCIe, gargalo crítico em
  arquiteturas híbridas.
\item
  \textbf{FullGPU é a campeã de qualidade:} Com um gap médio de
  $\mathbf{0,88\% \pm 1,14}$, a versão totalmente residente em GPU supera
  as variantes híbridas (1,28\% e 1,30\%, respectivamente). Uma explicação
  direta é que a variante FullGPU foi executada nos \emph{benchmarks}
  originais com 50 iterações de busca $2$-opt por indivíduo, contra apenas
  10 nas variantes híbridas e CPU (ver TODOs Pendentes). Essa assimetria
  favorece a convergência para melhores ótimos locais. O fato
  de que a variante Híbrida Otimizada — que também elimina transferências
  repetidas via processamento em lote — não atinja a mesma qualidade
  corrobora a hipótese de que o número superior de iterações é o fator
  determinante, e não apenas a permanência dos dados na GPU.
\end{enumerate}

Para validar estatisticamente as diferenças de qualidade (Gap) entre as
variantes implementadas em GPU (para todas as 38 instâncias), aplicou-se
o teste não-paramétrico de Friedman. O teste rejeitou a hipótese nula
com alta confiança (\(p \approx 6,9 \times 10^{-5}\)), confirmando que
há diferença global significativa no ranqueamento de desempenho. Em
seguida, o pós-teste de Nemenyi determinou pareadamente quais diferenças
são reais: a variante FullGPU mostrou-se estatisticamente superior à
HybridNaive (\(p \approx 0,00028\)) e à HybridOptimized
(\(p \approx 0,00071\)). Contudo, não houve diferença estatística
significativa entre as duas versões híbridas (\(p \approx 0,97\)). Isso
atesta que o incremento de qualidade da abordagem puramente residente
não é estocástico, mas sim estrutural.

\subsection{4.4 Melhor Algoritmo por Categoria de
Tamanho}\label{sec-44-melhor-algoritmo}

A Tabela\textasciitilde{}\ref{tbl:best-algorithm} segmenta o desempenho
``vencedor'' por faixa de tamanho da instância, considerando o critério
de menor \emph{gap} (qualidade).

\begin{longtable}[]{@{}llc@{}}
\caption{Melhor Algoritmo (Menor Gap) por Tamanho
\label{tbl:best-algorithm}}\tabularnewline
\toprule\noalign{}
Categoria & Melhor Algoritmo & Gap Médio \\
\midrule\noalign{}
\endfirsthead
\toprule\noalign{}
Categoria & Melhor Algoritmo & Gap Médio \\
\midrule\noalign{}
\endhead
\bottomrule\noalign{}
\endlastfoot
\textbf{Pequeno (\(n_{coords} \le 100\))} & \textbf{FullGPU} & 0,23\% \\
\textbf{Médio (\(100 < n_{coords} \le 400\))} & \textbf{FullGPU} &
0,79\% \\
\textbf{Grande (\(n_{coords} > 400\))} & \textbf{FullGPU} & 2,83\% \\
\end{longtable}

A dominância da variante \textbf{FullGPU} em todas as categorias de
tamanho reforça a observação de que a arquitetura residente é superior
para a qualidade da solução final. Mesmo que a variante
\textbf{HybridOptimized} seja mais rápida, ela tende a ficar presa em
ótimos locais ligeiramente piores com mais frequência. Isso sugere que,
para aplicações onde a qualidade da rota é crítica e o tempo de resposta
é secundário (mas ainda deve ser razoável), a abordagem FullGPU é a
recomendada. Para cenários de tempo real ou onde o \emph{throughput} é
prioritário, a HybridOptimized é a escolha ideal.

% TODO (Ação 1.1): CONCLUSÃO INVALIDADA — A "dominância da FullGPU em todas as
% categorias" pode ser artefato direto da assimetria paramétrica (50 vs 10 iterações
% de 2-opt). Após re-execução com parâmetros normalizados (Ação 3.1), esta seção
% deve ser reavaliada. A superioridade pode se manter, diminuir ou inverter.

\chapter{Discussão}\label{sec:discussao}

Os resultados validados no Capítulo 4 confirmam a hipótese de que o uso de GPUs acelera significativamente meta-heurísticas para o TSP. Contudo, demonstram que "usar GPU" não é uma solução trivial ou monolítica: a forma como a paralelização e o tráfego de dados são arquitetados define o perfil de desempenho de maneira categórica.

A abordagem \textbf{Híbrida Ingênua (Naive)}, embora funcional, sofre severamente com o gargalo de comunicação PCIe, servindo essencialmente como um modelo base metodológico. A abordagem \textbf{Híbrida Otimizada (Batch)}, ao processar avaliações em lotes, minimizou o custo de transferência e entregou a maior velocidade bruta, sendo ideal para aplicações práticas que exijam resposta rápida. Por sua vez, a abordagem \textbf{Totalmente Residente (FullGPU)} obteve o melhor refinamento médio das soluções (menor \emph{gap}), indicando que a eliminação total do \emph{overhead} de comunicação permite que a GPU concentre seus recursos na efetiva exploração do espaço de busca. A inviabilidade da execução da CPU para instâncias $n_{coords} > 100$ no desenho experimental atual reforça a utilidade da aceleração em hardware para a resolução de problemas de roteamento de médio e grande porte.

\section{Limitações da Análise de Desempenho}\label{sec-51-limitacoes-profiling}

A presente análise de desempenho baseia-se exclusivamente em medições macroscópicas de tempo total de execução, sem utilizar ferramentas de \emph{profiling} de hardware como o \emph{NVIDIA Nsight Compute}. Tais ferramentas, embora capazes de fornecer contadores detalhados por kernel (ocupação, pressão de registradores, divergência de \emph{warps}), serializam a execução dos kernels e exigem múltiplos passes de repetição, o que distorce substancialmente o tempo real em algoritmos iterativos de longa duração como o GA aqui implementado \autocite{nvidia2024cuda}.

Consequentemente, as hipóteses levantadas sobre causas arquiteturais no Capítulo~4 permanecem como conjecturas a serem validadas. Uma abordagem menos intrusiva, como o uso de temporizadores via primitivas CUDA na camada Python/CuPy, será avaliada como alternativa prática de instrumentação.

\section{Comparação com o Estado da Arte}\label{sec-52-comparacao-estado-arte}

Embora as abordagens paralelas desenvolvidas tenham alcançado alta velocidade de execução, é importante situar as limitações do GA+$2$-opt implementado frente ao estado da arte contemporâneo no TSP. Algoritmos dedicados e altamente especializados, notadamente o algoritmo LKH-3 \autocite{helsgaun2000effective} e o algoritmo genético baseado em Edge Assembly Crossover (EAX) \autocite{nagata2013powerful}, rotineiramente atingem $0\%$ de \emph{gap} (garantindo a solução ótima) na esmagadora maioria das instâncias padrão da TSPLIB. 

O GA+$2$-opt aqui explorado adota operadores canônicos e uma arquitetura base focada na investigação didática das relações de custo-benefício da programação paralela. Consequentemente, este trabalho foca na comparação entre as variantes do mesmo algoritmo (analisando os ganhos estruturais de cada estratégia de implantação), e não na competição isolada por poder de refinamento contra as meta-heurísticas de última geração que utilizam buscas como $5$-opt ou estratégias dinâmicas complexas de vizinhança. Uma investigação futura poderia avaliar a integração de operadores mais sofisticados (como o LK-move) às implementações em GPU, visando estreitar essa lacuna de qualidade.

\section*{TODOs Pendentes}

\begin{itemize}
    \item \textbf{Injustiça Paramétrica (Limite de Iterações no 2-opt):} Foi identificado que a variante \texttt{FULLGPU} foi executada nos experimentos originais com 50 passos rígidos na busca local, contra apenas 10 iterações na versão base da CPU e das abordagens híbridas. O parâmetro \texttt{two\_opt\_iterations} já foi ajustado no código (\texttt{ga\_fujimoto.cu}), sendo agora imprescindível a re-execução dos experimentos para obter um ranqueamento de lacunas (ìndice \emph{gap}) estritamente iso-algorítmico.
    \item \textbf{Análise Gráfica da Variância:} Investigar graficamente a distribuição da variância e dos valores atípicos da lacuna reportados na Tabela 2, com enfoque especial na assimetria e desvios padrão que se apresentaram matematicamente anômalos.
    \item \textbf{Instrumentação via Eventos CUDA:} Implementar temporizadores não-intrusivos via primitivas de eventos CUDA na camada Python/CuPy para medir a duração individual de cada núcleo de execução sem a distorção causada pelo \emph{Nsight Compute} \autocite{nvidia2024cuda}. Esta é uma etapa prioritária para identificar os gargalos reais na GPU antes de qualquer outra investigação arquitetural.
    \item \textbf{Investigação de Pressão de Registradores:} Validar, via \emph{NVIDIA Nsight Compute}, se a alta demanda de registradores pelos núcleos de execução do algoritmo memético (especialmente o cruzamento OX e a busca $2$-opt) reduz a ocupação dos Multiprocessadores de Fluxo (SMs) e contribui para a degradação de velocidade em instâncias grandes.
    \item \textbf{Investigação de Divergência de Grupos de Execução:} Avaliar empiricamente se a natureza sequencial do operador de cruzamento OX --- com laços condicionais dependentes de dados --- causa serialização excessiva dentro dos grupos de 32 linhas de execução (\emph{warps}), impactando a vazão efetiva.
    \item \textbf{Análise de Complexidade Algorítmica vs.\ Arquitetural:} Separar formalmente a contribuição de limitações do algoritmo memético escolhido --- incluindo todas as suas fases: geração da população inicial, seleção, cruzamento (OX), mutação e busca local ($2$-opt) --- do impacto de limitações de \emph{hardware} da GPU. Considerar também a exploração de variantes algorítmicas alternativas (outros operadores de cruzamento, mutação adaptativa, buscas locais mais agressivas) que possam ser intrinsecamente mais adequadas à arquitetura paralela massiva.
\end{itemize}

\chapter{Conclusões e Trabalhos Futuros}\label{cap6-conclusoes}

O presente relatório documenta o progresso na investigação e quantificação dos benefícios de portar um Algoritmo Memético (GA+2-opt) para GPU visando resolver o Problema do Caixeiro Viajante. O desenvolvimento sistemático de quatro variantes (CPU, Híbrida Ingênua, Híbrida Otimizada em lotes e Totalmente Residente na GPU) confirmou que a placa de vídeo acelera o processo decisivamente, contanto que o gargalo de comunicação com a placa mãe seja evitado. Os resultados atuais refletem as escolhas específicas de cada fase do algoritmo memético --- geração da população, seleção, cruzamento (OX), mutação e busca local ($2$-opt) --- e a investigação de variantes alternativas para cada uma dessas fases constitui trabalho futuro relevante.

Os resultados empíricos validados sobre as 38 instâncias da biblioteca TSPLIB demonstraram que a minimização da transferência de dados é o fator primordial de sucesso. A variante \textbf{GA-Híbrido-Otimizado} destacou-se pela velocidade (ganho mediano de até 5,69$\times$ e milhares de vezes mais rápida que a execução sequencial), sendo a abordagem recomendada quando a urgência da resposta e a vazão de processamento são prioritárias. Por outro lado, a versão \textbf{GA-FullGPU} garantiu o melhor refinamento médio de rotas ($0,88\%$ de gap), comprovando que alocar todos os dados na VRAM da placa de vídeo evita interrupções disruptivas causadas pelo barramento PCIe. 

\section{Trabalhos Futuros}\label{sec-61-trabalhos-futuros}

O aprofundamento das investigações da arquitetura CUDA para problemas de roteamento apresenta caminhos lógicos imediatos. Como extensões práticas deste trabalho, sugerem-se:

\begin{itemize}
    \item \textbf{Implementação de Ferramentas de Perfilação Adequadas:} Como o monitoramento com o \emph{NVIDIA Nsight Compute} pode ser excessivamente intrusivo \autocite{nvidia2024cuda}, propõe-se avaliar contornos práticos de instrumentação de \emph{software} (como temporizadores via eventos CUDA na camada Python) para cronometrar a nível de núcleo de execução e isolar falhas de ocupação sem distorcer o tempo real de execução \autocite{kirk2010programming}.
    \item \textbf{Aprimoramento via Primitivas de Comunicação Intra-grupo:} Refatorar o operador de cruzamento ordenado (OX), que lida de forma imperativa com a construção vetorial, para aplicar as chamadas de comunicação direta entre linhas de execução ($\operatorname{__shfl\_sync}$) \autocite{nvidia2024cuda}. Essa abordagem tem o potencial de reduzir a divergência de execução dentro do grupo, pois compartilharia o dado via registradores diretamente entre as linhas de execução \autocite{harris2007optimizing}.
    \item \textbf{Padronização para Aritmética de Inteiros Discretos:} Substituir todos os cálculos geométricos euclidianos que dependem de ponto flutuante contínuo (\emph{float64}) por rotinas de cálculo com inteiros (\emph{int64}), alinhando estritamente a matemática das distâncias à função de arredondamento (\texttt{nint}) definida na documentação canônica da TSPLIB \autocite{reinelt1991tsplib}. Isso extinguirá pequenas flutuações e erros de aproximação causados pela aceleração vetorial em GPU.
    \item \textbf{Exploração de Variantes Algorítmicas:} Investigar operadores alternativos para cada fase do algoritmo memético --- cruzamento (PMX, CX), mutação adaptativa, buscas locais mais agressivas (Or-opt, $3$-opt) --- visando identificar combinações intrinsecamente mais adequadas à arquitetura paralela massiva da GPU \autocite{larranaga1999genetic, fujimoto2011highly}.
\end{itemize}

\printbibliography

\end{document}
